\documentclass[11pt]{article}

\usepackage[T1]{fontenc}
\usepackage[utf8]{inputenc}
\usepackage[english]{babel}
\usepackage[margin=1in]{geometry}
\usepackage{lmodern}
\usepackage{microtype}
\usepackage{mathtools}
\usepackage{amsmath,amssymb,amsthm}
\usepackage{array}
\usepackage{booktabs}
\usepackage{tabularx}
\usepackage{float}
\usepackage{graphicx}
\usepackage{pdflscape}
\usepackage{caption}
\usepackage{enumitem}
\usepackage{fancyhdr}
\usepackage[numbers]{natbib}
\usepackage{doi}
\usepackage{xurl}
\usepackage{hyperref}
\usepackage{cleveref}
\usepackage{orcidlink}
\usepackage{titling}



\newcommand{\gradeS}{\textsf{[S]}}
\newcommand{\gradeR}{\textsf{[R]}}
\newcommand{\gradeO}{\textsf{[O]}}
\newcommand{\gradeOpen}{\textsf{[OPEN]}}

\newcommand{\Kviab}{K_{\mathrm{viab}}}
\newcommand{\Kbw}{K_{\mathrm{bw}}}
\newcommand{\Rrec}{R_{\mathrm{rec}}}
\newcommand{\Cpkg}{\mathfrak{C}}
\newcommand{\Certsimp}{\mathrm{Cert}_{\mathrm{simp}}}
\newcommand{\Dfact}{\Delta_{\mathrm{fact}}}
 
\newtheorem{theorem}{Theorem}[section]
\newtheorem{lemma}[theorem]{Lemma}
\newtheorem{proposition}[theorem]{Proposition}
\newtheorem{corollary}[theorem]{Corollary}
\theoremstyle{definition}
\newtheorem{definition}[theorem]{Definition}
\newtheorem{remark}[theorem]{Remark}
\newtheorem{construction}[theorem]{Construction}
\newtheorem{postulate}{Postulate}
\floatstyle{ruled}
\newfloat{algorithm}{tbp}{loa}
\floatname{algorithm}{Algorithm}

\urlstyle{same}
\captionsetup{
  font=small,
  labelfont=bf,
  labelsep=period
}
\captionsetup[table]{skip=6pt,position=top}
\captionsetup[figure]{skip=6pt}
\captionsetup[algorithm]{skip=6pt,position=top}
\crefname{algorithm}{algorithm}{algorithms}
\Crefname{algorithm}{Algorithm}{Algorithms}
\crefname{theorem}{theorem}{theorems}
\Crefname{theorem}{Theorem}{Theorems}
\crefname{lemma}{lemma}{lemmas}
\Crefname{lemma}{Lemma}{Lemmas}
\crefname{proposition}{proposition}{propositions}
\Crefname{proposition}{Proposition}{Propositions}
\crefname{corollary}{corollary}{corollaries}
\Crefname{corollary}{Corollary}{Corollaries}
\crefname{definition}{definition}{definitions}
\Crefname{definition}{Definition}{Definitions}
\crefname{remark}{remark}{remarks}
\Crefname{remark}{Remark}{Remarks}
\setlist[enumerate]{leftmargin=*, itemsep=2pt, topsep=4pt}
\setlength{\droptitle}{-3.2em}
\pretitle{\begin{center}\LARGE\bfseries}
\posttitle{\par\end{center}\vspace{0.5em}}
\preauthor{\begin{center}\normalsize}
\postauthor{\par\end{center}\vspace{-0.4em}}
\predate{\begin{center}\small}
\postdate{\par\end{center}\vspace{0.6em}}
\setlength{\emergencystretch}{2em}
\clubpenalty=10000
\widowpenalty=10000
\displaywidowpenalty=10000
\interfootnotelinepenalty=10000
\allowdisplaybreaks[2]
\fancypagestyle{firstpage}{\fancyhf{}
  \fancyfoot[C]{\scriptsize Automorph Inc., Wilmington, DE, USA \quad \texttt{ioannis@automorph.io} \quad \textcopyright\ 2026 \quad CC-BY 4.0}
  \renewcommand{\headrulewidth}{0pt}
  \renewcommand{\footrulewidth}{0pt}
}

\title{An Audited Closure Architecture for Consciousness:\\
Exact Finite Models, Structural Theorems, and Conditional Realization}
\author{Ioannis Tsiokos\,\orcidlink{0009-0009-7659-5964}}
\date{\today}
\hypersetup{
  pdftitle={An Audited Closure Architecture for Consciousness: Exact Finite Models, Structural Theorems, and Conditional Realization},
  pdfauthor={Ioannis Tsiokos},
  hidelinks
}

\begin{document}

\maketitle
\thispagestyle{firstpage}

\begin{abstract}
Formal theories of consciousness often mix mathematical structure, empirical
access, and ontological interpretation inside single claims. We present an
architecture that separates them by construction: an eight-clause closure
package (objecthood, strict promotion, a reflexive self-coordinate,
non-exposure, action credit, temporal closure, binding, and viability)
specified as a full design and audit contract, with every result graded as a
structural theorem, a recognition bridge, a conditional ontological
implication, or an open target. The evidence base consists of exact finite
studies, meaning exhaustively enumerated rational computations on declared
carriers, together with selected machine-checked Lean kernels. The central exact
result: on a declared twelve-state deadline carrier, every public-lens policy
has an empty viability kernel, the canonical predictive quotient restores a
ten-state kernel, and the unique minimal closure-restoring refinement splits
exactly one public fiber; extension beyond the carrier is a conditional
schema, not a generic theorem. A structural catalog then recasts privacy, the
explanatory gap, zombies, free will, binding, temporal experience, valence,
quality geometry, inverted spectra, personal identity, combination, and the
meta-problem as interface-relative factorization, quotient, viability, gauge,
and audit results. Realization is split into seven graded claims: a mechanized
conditional bridge purchases it only from assumed certificate soundness plus
an ontic-closure postulate, and there is no integrated certified instance and
no empirical validation. An auditable code baseline with executable
finite controls and predeclared prospective empirical obligations accompanies
the paper.
\end{abstract}

\medskip
\noindent\textbf{Keywords.} mathematical consciousness science; structural
theory of consciousness; exact finite-state models; predictive-state quotient;
viability kernel; observational equivalence; formal verification.

\medskip
\noindent\textbf{Artifact note.} Every value in a numerical table in this
paper is read from a deterministic, schema-validated snapshot
(\texttt{paper/data/current\_results.json}) exported from the repository's
Python compute objects; the remaining in-text exact values are pinned by the
repository's theorem catalog and test suites, which the same artifact
contains. An immutable public artifact
identifier (DOI and frozen commit) will be added at release; the local
verification baseline is reported in Appendix~\ref{app:repro}.

\section{Problem, result classes, and contribution map}
\label{sec:intro}


Formal theories of consciousness are difficult to compare, and difficult to
trust, partly because three different kinds of claim routinely travel in one
sentence: mathematical claims about a formal structure, empirical claims about
how that structure meets observation, and ontological claims about what the
structure's satisfaction would \emph{be}. The field has begun to treat this as
a foundational problem rather than a stylistic one. General frameworks for
mathematical models of consciousness ask what a model must declare before it
says anything \citep{Kleiner2020MathematicalModels, KleinerTull2021}; formal
work on testability shows that inference procedures and prediction data can
interact so badly that a theory becomes unfalsifiable or trivially falsified
\citep{KleinerHoel2021, DoerigEtAl2020HardCriteria}; and large adversarial
collaborations have demonstrated how much machinery is needed to test even two
well-resourced theories against each other, without producing a definitive
verdict \citep{CogitateConsortium2025}. Landscape reviews count theory families
by the dozen \citep{SethBayne2022}. In such a field, a paper's first obligation
is to say exactly what kind of claim each of its sentences makes.

This paper takes that obligation as its organizing principle. It presents an
architecture for consciousness built from standard finite mathematics:
finite carriers, lenses, quotients, refinements, closure operators, viability
kernels, and audited promotion between description layers. Around that
architecture it assembles a catalog of exact finite results, a small set of
machine-checked kernels, and one explicitly priced ontological bridge. Every
scientific sentence in the paper is assigned to one of four \emph{claim lanes},
defined in Section~\ref{sec:claims}: structural results \gradeS{} (theorems and
exact finite constructions), recognition bridges \gradeR{} (proposed
interpretive mappings, never derivations), ontological implications \gradeO{}
(conditional on named postulates), and open targets \gradeOpen. The lane
discipline is not decoration; it is enforced by a sentence-level claim ledger
in the artifact, and the paper's central negative rule is that no amount of
\gradeS{} material ever implies an \gradeO{} conclusion without the named
purchases.

\subsection{Research questions}
\label{sec:questions}

The paper is organized around five questions.

\begin{enumerate}
\item Under what exact finite conditions does viability under a bounded public
self-lens force a predictive self-coordinate, or a closure-restoring refinement
of the public description? (Section~\ref{sec:t1}.)
\item What component results and conditional dependencies follow when
objecthood, strict promotion, reflexivity, hiddenness, drive, temporal closure,
binding, and viability are specified by one audit contract, while the
current witnesses generally live on separately declared carriers?
(Sections~\ref{sec:arch} and~\ref{sec:clusters}.)
\item Which classical consciousness problems (privacy, the explanatory gap,
zombies, free will, binding, the specious present, inverted spectra, personal
identity, combination, the meta-problem) become interface-relative
factorization, quotient, viability, gauge, audit, or transport questions under
that contract? (Section~\ref{sec:clusters}.)
\item Which results are exact finite constructions, which extend to conditional
schemas, which are recognition bridges, and which remain open?
(Sections~\ref{sec:claims}--\ref{sec:t17} throughout.)
\item What ontological interpretation follows, and only follows, after
explicitly purchasing certificate soundness and Ontic Closure Completeness?
(Section~\ref{sec:t17}.)
\end{enumerate}

\subsection{Evidence base and method}

The evidence base is deliberately narrow and fully auditable. Every law in the
catalog is landed as an \emph{exact finite study}: a small, fully declared
finite carrier on which the relevant property is computed by exhaustive
enumeration in exact rational arithmetic, together with declared controls that
show what removes the effect. These are computational witnesses and exhaustive
enumerations, not behavioral or empirical experiments; no empirical data of any
kind were collected, fitted, or used to validate the architecture
(Section~\ref{sec:empirical}). Selected results are additionally mechanized:
each of the nineteen law modules has a Lean kernel, and the exact scope
boundary of each kernel (what Lean proves versus what it pins as certified
constants matching the Python enumeration) is stated where the result is used
and cataloged in Appendix~\ref{app:repro}. ``Machine-checked'' in this paper always
modifies a named declaration, never the catalog as a whole.

The architecture is built on an inter-layer grammar developed in a series of
companion papers (the ``Six Birds'' corpus), which supplies audited notions of
promotion between description layers, factorization defects, no-overreading
rules for claims that cross layers, and limits on same-level self-audit. That
corpus is background machinery with its own hypotheses, cited exactly where a
theorem is imported (Appendix~\ref{app:imports}); it is not independent
evidence for anything about consciousness, and the paper is written to be read
without prior exposure to it: every load-bearing definition is given first in
standard mathematical language (Section~\ref{sec:setting}).

\subsection{The architecture and the two headline boundaries}

Section~\ref{sec:arch} defines the central object: an eight-clause closure
package, the \emph{C-package} $\Cpkg$. It specifies what a candidate
structure must exhibit: multiscale objecthood (C1), strict promotion over a
public layer (C2), a recursively stable self-coordinate (C3), non-exposure of
the determining state at the declared interface (C4), audited self/other action
credit (C5), a funded temporal order/tick/arrow structure (C6), strict
recombination binding (C7), and nonempty maintained viability (C8), together
with cross-clause coherence conditions and an audit-gate panel. The package is
a declared \emph{target}, not an exhibited system.

Two results bound everything else in the paper.

First, the strongest exact result (Section~\ref{sec:t1}): on a fully declared
twelve-state deadline carrier, \emph{every} public-lens feedback policy has an
empty viability kernel (by exhaustive enumeration, not sampling), while the
canonical predictive quotient restores a ten-state kernel; and among all $64$
refinements of the public lens, the \emph{unique minimal} closure-restoring
refinement is a seven-class quotient that splits exactly one public fiber.
Under a deadline, on this carrier, survival forces a predictive
self-coordinate. The result is theorem-grade on its declared carrier; its
family-level extension is a conditional schema whose hypotheses are displayed,
not a generic theorem.

Second, the visibly open boundary (Section~\ref{sec:t17}): the realization
claim, that satisfying the certificate \emph{is being} a subject rather than
imitating one, is split into seven separately graded objects. A simplified
certificate schema and a conditional bridge theorem are mechanized in Lean.
The bridge's conclusion is purchased by two premises that are assumed, not
proved: certificate soundness, and the postulate that lawful audited closure
is realization. The integrated instance (one carrier satisfying a full
certificate with a \emph{proved} soundness theorem) is recorded as open, with
an explicit \texttt{false} marker in the artifact. A reader who declines the
premises keeps the entire structural catalog and loses only the ontological
reading. The refusal is priced and respected.

\subsection{Contributions}
\label{sec:contributions}

In order of importance:

\begin{enumerate}
\item \textbf{A claim architecture.} An explicit, artifact-enforced separation
among Tier~S structural results, Tier~R recognition bridges, Tier~O conditional
interpretation, and open targets, applied to every sentence of a consciousness
theory (Section~\ref{sec:claims}).
\item \textbf{A declared architecture and audit contract.} The eight-clause
C-package as a full design target with component and interaction audits, kept
explicitly distinct from the simplified certificate schema that is currently
mechanized, and from the still-open integrated instance
(Sections~\ref{sec:arch},~\ref{sec:t17}).
\item \textbf{An exact forcing result.} The deadline carrier on which public
self-description provably fails and a unique minimal predictive refinement
restores viability, with exhaustive policy and refinement enumeration, closure
and perturbation audits, and controls (Section~\ref{sec:t1}).
\item \textbf{A structural catalog.} Exact, interface-relative component
results recasting privacy, the explanatory gap, zombies, reports, self-audit,
binding, temporal closure, constraint prices, quality geometry, gauge freedom,
lineage, combination, and the meta-problem as factorization, quotient,
viability, gauge, audit, and transport facts on declared finite carriers
(Section~\ref{sec:clusters}).
\item \textbf{A graded realization program.} The seven-way T17 split with
executable negative attacks, an axiom-relative conditional bridge, and an
explicit open integrated target (Section~\ref{sec:t17}).
\item \textbf{An audited artifact baseline.} Exact rational Python studies with
191 passing finite-study tests, a deterministic schema-validated nineteen-law
result snapshot with 8 passing snapshot tests, nineteen Lean law modules in a
36-job build, dependency and contract auditors, and a declared inherited trust
base, with release-readiness gaps reported explicitly
(Appendix~\ref{app:repro}).
\end{enumerate}

Empirical status is a limitation and a future obligation, not a seventh
contribution: no empirical validation exists, an inherited literature battery
failed its own admission rules and is retained only as non-evidential protocol
history, and the prospective empirical interface is specified as obligations
(Section~\ref{sec:empirical}).

We do not claim priority for mathematical consciousness science, axiomatic
models, privacy or unreadability results, unity, gauge symmetry, quality
spaces, predictive-state quotients, viability kernels, shadow prices, or
self-reference limits taken separately; each has its own literature, credited
in Section~\ref{sec:related}. The defensible novelty is the particular audited
composition: the claim discipline, the exact carriers, the forcing result, and
the conditional realization split.

\subsection{What this paper does not do}

Five nonclaims frame everything that follows. (i)~No empirical experiment is
reported and no empirical prediction is tested. (ii)~No candidate subject is
proposed, examined, or certified; no actual system is claimed to be conscious
or not conscious. (iii)~No full classifier for consciousness is offered; the
certificate schema that exists is simplified, and its integrated instantiation
is open. (iv)~No unconditional ontological conclusion is drawn anywhere; Tier~O
content is confined to one conditional bridge with named premises.
(v)~No claim is made to be the first mathematical theory of consciousness; the
paper works inside an existing field
\citep{Kleiner2020MathematicalModels, SignorelliWangCoecke2021,
ChangEtAl2020InformationClosure, KleinerTull2021}.

\subsection{Reader's map}

Section~\ref{sec:claims} defines the claim lanes and the minimal formal
setting. Section~\ref{sec:arch} defines the C-package.
Section~\ref{sec:t1} develops the exact forcing result.
Section~\ref{sec:clusters} organizes the remaining laws into four consequence
clusters: interface and access; self-audit and action credit; the structure of
experience; persistence and boundaries. Section~\ref{sec:t17} prices the
realization bridge and displays the open target. Section~\ref{sec:empirical}
states the empirical status. Section~\ref{sec:related} positions the work.
Section~\ref{sec:limits} collects limitations and concludes. Appendices give
imported machinery (\ref{app:imports}), full T1 details (\ref{app:t1}), the
per-law result compendium (\ref{app:compendium}), the T17 schema and attack
corpus (\ref{app:t17}), the non-evidential empirical-interface history
(\ref{app:empirical}), reproducibility and formalization (\ref{app:repro}),
and the claim and overlap ledgers (\ref{app:ledgers}).

 \section{Claim types and the minimal formal setting}
\label{sec:claims}


This section fixes the two pieces of infrastructure that everything else
depends on: a typing discipline for claims, and the small stock of standard
finite mathematics in which all results are stated.

\subsection{Four claim lanes}
\label{sec:lanes}

Every scientific sentence in this paper belongs to exactly one lane.

\begin{description}
\item[Tier S (structural).] Mathematical content about declared finite
objects. Tier~S subdivides by \emph{claim form}: a \emph{theorem-grade exact
result} is proved or exhaustively computed on its declared carrier; a
\emph{schema} is a definition or conditional statement whose hypotheses are
displayed and not discharged; a \emph{pointer} re-scopes an existing result to
a new reading without new mathematics; a \emph{composition} combines landed
results under a declared convention. Tier~S claims are established relative to
their carriers and interfaces, and to nothing else.
\item[Tier R (recognition).] A proposed mapping between a structural object
and a familiar mental or phenomenal category, for example reading a
constraint's dual price as felt valence. Recognition bridges are hypotheses.
They are \emph{proposed}, \emph{named}, and \emph{retractable}; they are never
derivations, and a counterexample to a bridge retracts the bridge, not the
mathematics under it.
\item[Tier O (ontological).] A claim about what satisfaction of the structure
would be. All Tier~O content in this paper is conditional on named postulates
(Section~\ref{sec:t17}); the paper contains no unconditional Tier~O sentence.
\item[OPEN.] A target that is stated, motivated, and not achieved. Open
targets are recorded in the artifact with explicit markers rather than
implied by silence.
\end{description}

Table~\ref{tbl:grades} summarizes the lanes, their evidential support, and the
inference each lane blocks. Figure~\ref{fig:lanes} shows the same discipline as
a flow constraint: computational and literature inputs enter only their proper
lanes, and no arrow runs from computation or literature directly to
realization.

\begin{table}[t]
\centering\small
\begin{tabular}{p{0.1\linewidth}p{0.2\linewidth}p{0.22\linewidth}p{0.18\linewidth}p{0.2\linewidth}}
\toprule
Grade & Object & Evidential support & Licensed wording & Blocked inference \\
\midrule
\gradeS{} theorem/ exact & declared finite carrier and result & proof, or exhaustive exact enumeration with tests; selected Lean kernels & ``proves on the declared carrier,'' ``computes exactly,'' ``exhaustively enumerates'' & generalization beyond the carrier; any empirical or phenomenal conclusion \\
\gradeS{} schema & conditional statement with displayed hypotheses & definition-level checking; conditional proofs & ``establishes under assumptions,'' ``gives a conditional schema'' & unconditional use; silent discharge of hypotheses \\
\gradeR{} & proposed structure-to-phenomenon mapping & argued fit of profile; no derivation & ``proposes,'' ``recognizes,'' ``interprets, conditional on'' & treating the mapping as derived or measured \\
\gradeO{} & conditional ontological implication & named postulates plus a mechanized conditional theorem & ``implies after assuming,'' ``is conditional on,'' ``purchases'' & detaching the conclusion from its premises \\
\gradeOpen{} & stated, unachieved target & explicit negative markers in the artifact & ``remains open,'' ``is not exhibited'' & ``follows in principle,'' ``is nearly proved'' \\
\bottomrule
\end{tabular}
\caption{\textbf{Tbl-1: claim-grade legend.} Each grade names its object, its
support, its licensed verbs, and the inference it blocks. Scope/price: the
legend is a reporting discipline for this paper, not a proposed norm for the
field.}
\label{tbl:grades}
\end{table}

\begin{figure}[t]
\centering
\fbox{\parbox{0.92\linewidth}{\small\emph{[Figure F1: claim and evidence
lanes; rendered per the artifact figure plan.]} Four horizontal lanes:
structural result, recognition bridge, ontological implication, open
obligation. Python enumeration and Lean declarations feed only the structural
lane; philosophical postulates feed only the ontological lane; literature
feeds comparison, never evidence. No arrow connects computation or literature
directly to realization.}}
\caption{\textbf{F1: claim and evidence lanes.} Takeaway: every sentence has
one lane, and evidence types cannot cross lanes. Scope restriction: the
diagram is schematic; the sentence-level assignment lives in the artifact's
claim ledger.}
\label{fig:lanes}
\end{figure}

\subsection{The minimal formal setting}
\label{sec:setting}

All objects are finite. Probabilities and costs are exact rationals; no
floating-point value appears in any result.

\begin{definition}[Carrier, actions, lens]
A \emph{carrier} is a finite set $X$ of states. A \emph{transition/action
family} over $X$ is a finite set $A$ of actions with, for each $a\in A$,
either a map $F_a : X \to X$ or a row-stochastic kernel $P_a$ on $X$; whether
dynamics are deterministic or stochastic is declared per construction. A
\emph{lens} (observation map) is a map $q : X \to Y$ onto a finite value set;
its fibers form the \emph{observational partition} $\Pi_q$ of $X$. A
\emph{policy for $q$} is a feedback rule that selects actions measurably with
respect to $\Pi_q$ (more generally, with respect to the history of $q$-values).
\end{definition}

\begin{definition}[Quotient, refinement, factorization defect]
A partition $\Pi'$ \emph{refines} $\Pi$ when every block of $\Pi'$ is contained
in a block of $\Pi$. Given two lenses $\pi_0, \pi_1$ on the same carrier, the
\emph{factorization defect} is
\[
\Dfact(\pi_0,\pi_1) \;=\; \{(s,s') : \pi_0(s)=\pi_0(s') \text{ and }
\pi_1(s)\neq\pi_1(s')\},
\]
the set of \emph{split pairs}: pairs identified by the coarse description and
separated by the fine one. $\pi_1$ factors through $\pi_0$ (there is $\varphi$
with $\pi_1=\varphi\circ\pi_0$) if and only if $\Dfact(\pi_0,\pi_1)=\emptyset$.
A nonempty defect is the exact, checkable meaning of ``the fine content is not
recoverable from the coarse description'' used throughout this paper.
\end{definition}

\begin{definition}[Predictive and observational equivalence]
Fix a carrier with declared dynamics and a lens $q$. Two states (or histories)
are \emph{observationally equivalent} when all admissible current observables
agree on them; the induced quotient is the \emph{current public quotient} $Q$.
They are \emph{predictively equivalent} when they agree on the full family of
admissible \emph{future} behavior (continuation laws, or distributions over
future observable trajectories under all admissible continuations); the induced
quotient is the \emph{predictive quotient} $M$. Predictive equivalence refines
observational equivalence, giving a canonical comparison map $\pi : M \to Q$.
The gap between them, same $Q$-class with different $M$-class, is a
factorization defect and is the recurring engine of this paper.
\end{definition}

Predictive quotients are standard mathematics: they are the finite controlled
analog of causal-state equivalence in computational mechanics
\citep{CrutchfieldYoung1989, ShaliziCrutchfield2001}, and we import the
construction, not the stationary-process theorems (the distinction matters in
Section~\ref{sec:t1}).

\begin{definition}[Closure operator and audit]
A \emph{closure operator} on a finite structured set is an extensive, monotone,
idempotent endomap; its fixed points are the descriptions the structure
recognizes as complete. When a closure-like operator $E$ is only approximately
idempotent, its \emph{idempotence defect} $\delta(E)=\tfrac12\max_z
\lVert E^2(z,\cdot)-E(z,\cdot)\rVert_1$ is computed exactly and reported rather
than assumed small. An \emph{audit} is a declared, executable check on a
construction, such as a defect computation, a control run, or an exhaustive
comparison, whose outcome is recorded with the result. A \emph{gate} is a named audit
whose passage is required before a claim may be graded; the gate vocabulary
used here (sufficiency, descent, stability, control-of-content, no-smuggling,
visibility, audit, strictness) is imported from the interaction calculus of
the companion corpus (Appendix~\ref{app:imports}) and is used as a reporting
format, not re-proved.
\end{definition}

\begin{definition}[Viability kernel]
Fix a carrier $X$, a safety predicate, feasible action sets, and declared
successor supports. The \emph{viability kernel} $\Kviab$ is the greatest fixed
point of the controlled-invariance operator
\[
V(K) = \{\, s \in K : s \text{ safe and } \exists a \in A_{\mathrm{feas}}(s)
\text{ with } \mathrm{Post}(s,a) \subseteq K \,\},
\]
computed by descending iteration; this is the standard discrete-time viability
kernel \citep{AnayaEtAl2009Viability, AubinBayenSaintPierre2011}. For a lens
$\ell$, we write $\Kviab(\ell)$ for the best kernel achievable by
$\ell$-measurable policies. A nonempty kernel means controlled persistence in
the model, nothing more.
\end{definition}

\subsubsection*{The trust chain}

Every result travels the same four-step chain, and the paper reports which
steps each result completed. (1)~A \emph{mathematical statement} on a declared
carrier, with grade and nonclaims, in the theorem catalog. (2)~A \emph{Python
enumeration}: exact rational computation of every quantity, with declared
controls, under test. (3)~A \emph{Lean declaration}: a named kernel per law,
whose scope boundary (what is proved versus what is pinned as certified
constants matching the Python computation) is part of the result.
(4)~An \emph{interpretation}, which is where lanes R and O begin and where the
ledger's blocked-inference column applies. Prose never outranks an executable
artifact: the theorem catalog controls statements, the current Python
computations control values, the Lean sources control what is machine-checked.

\subsection{Six Birds vocabulary, second}
\label{sec:sixbirds}

The companion corpus organizes the same mathematics under a six-role grammar
(threshold, budget, route, staging, packaging, and audit; P1--P6) with a typed
calculus of promotion between description layers: a \emph{promotion bridge}
carries content from a lower to a higher layer; \emph{strict} promotion means
the higher description does not factor through the lower one
($\Dfact\neq\emptyset$); \emph{no-smuggling} and \emph{no-overreading} rules
police what may be used and claimed across layers; and \emph{same-level
self-audit} is provably never accepted, while level-shifted, partial audit
is. This paper uses those results as imported theorems with exact citations
(Appendix~\ref{app:imports}); the six roles never organize sections, are never
renamed as neurocognitive faculties, and impose no requirement that every role
be active in every carrier. Analogy with a sibling result is never treated as
theorem transport.

 \section{The audited closure architecture}
\label{sec:arch}


This section defines the paper's central object. It is a \emph{contract}: a
declaration of what a structure must exhibit, under what audits, for the
catalog's results to compose. It is not an exhibited system, and
Section~\ref{sec:t17} keeps the distance between the contract and any
instantiation explicit.

\subsection{The full C-package}
\label{sec:cpackage}

\begin{definition}[Full C-package target]
\label{def:cpackage}
The \emph{full C-package} $\Cpkg$ is the eight-clause full-design and audit
contract $\mathrm{C1}$--$\mathrm{C8}$, together with cross-clause coherence
data and a full promotion-gate panel. A candidate consists of a carrier $X$,
declared transition/action families, a public lens and exposure channels, a
predictive self-lens, branch structure, a maintenance ledger, and protocol
generators; the clauses constrain that data as follows.
\begin{description}
\item[C1 (objecthood).] A declared multiscale exposure-closure family with
robust fixed-point structure: closure operators $E_\tau$ at declared stagings
$\tau$, exact fixed-point and defect bookkeeping, and perturbation audits with
declared retention thresholds.
\item[C2 (strict promotion).] A strict promoted distinction across declared
layers: $\Dfact(\pi_0,\pi_1)\neq\emptyset$ for the declared lower and promoted
readouts, certified by a split pair.
\item[C3 (reflexive self-coordinate).] A recursively stable self-coordinate:
refinement of the public self-lens by dynamical profiles terminates at a
predictive quotient that carries the package's own continuation coordinate.
\item[C4 (non-exposure).] The determining predictive state is not exposed
relative to the declared public channel family: no admissible current
observable factors through it non-constantly on the realized support.
Hiddenness is channel-relative by construction.
\item[C5 (drive/credit).] A declared audit of action credit: generator
provenance routes computed directionality into \emph{self}, \emph{other}, or
\emph{none}, with scheduled and internalized controls that must not be
credited as self.
\item[C6 (temporal closure).] A maintained, funded order/tick/arrow structure:
a staged phase carrier that persists and progresses under budget, with the
arrow carried by the package's own audit channel rather than by external
scheduling.
\item[C7 (binding).] Strict recombination refinement over the branchwise
quotient: joint content over declared branches, witnessed by
$\Rrec \succ \Kbw$ (Section~\ref{sec:structure}).
\item[C8 (viability).] Nonempty maintained viability: $\Kviab\neq\emptyset$
under the declared maintenance ledger, with underfunding collapsing the
kernel.
\end{description}
\end{definition}

The full design object also names four \emph{derived} structures that are
consequences of the clauses rather than clauses themselves: D1, the dual-price
structure of binding C8 constraints (Section~\ref{sec:structure}); D2, the
descended quality geometry of the package's own states
(Section~\ref{sec:structure}); D3, a proposed scheduling/marking reading of
attention-like allocation, which is recognition-grade only (the present
repository implements no price-based attention scheduler); and D4, the declared
report/descent channel with its no-overreading boundary
(Section~\ref{sec:interface}).

Two premises sit \emph{outside} the structural boundary and are
stated here only so the reader can see where they will be spent
(Section~\ref{sec:t17}): A1 (Ontic Closure Completeness) and A2 (Substrate
Indifference) are philosophical postulates, not clauses, and no structural
result in Sections~\ref{sec:t1}--\ref{sec:clusters} depends on them.

\begin{table}[t]
\centering\footnotesize
\begin{tabular}{p{0.06\linewidth}p{0.2\linewidth}p{0.24\linewidth}p{0.22\linewidth}p{0.18\linewidth}}
\toprule
Clause & Standard object & Strongest landed evidence & Simplified (T17-S) proxy & Blocking condition / nonclaim \\
\midrule
C1 & closure operators $E_\tau$ with fixed-point, defect, and perturbation audits & T1: exact $E_\tau$ audit on the ten-state kernel; unique perturbed fixed point retaining $77/80$ CORE mass at $\varepsilon=1/16$ & mechanics-coupled closure with staged-transition derivation and one-step audits & a fixed point is not an object or phenomenal entity without a recognition bridge \\
C2 & factorization defect $\Dfact\neq\emptyset$ & T1 split pair; T3 nonfactorization corollary & whole-carrier defect from strict ledger & no experiential novelty or universal explanatory gap \\
C3 & terminal profile refinement to predictive quotient & T1: kernel-restricted refinement $6\to7\to7$, terminal $=$ canonical $M$ on the kernel & one-step complete-observation-stable quotient & one-step stability at generic grade is not recursive terminality in full generality \\
C4 & non-exposure of determining state & T2: connector loop moves $M$ while fixing $Q$; dissolution control & single collision-generated class under the complete joint observation & no absolute privacy; expanded interfaces may expose \\
C5 & credit rule over generator provenance & T5: exact ledgers credit self/other/none; DPI null blocked & aggregate no-fake-arrow balance control & aggregate balance does not instantiate the concrete two-component lift \\
C6 & funded order/tick/arrow & T8: funded trace with monotone ledger, $\Delta A = 2$; unfunded stall & budget-disciplined trace proxy & the trace proxy is not the full imported clock theorem \\
C7 & strict recombination refinement & T7: $\Rrec\succ\Kbw$ witness; marking/erasure controls & global refinement with nontrivial fiber & no phenomenal-unity derivation \\
C8 & viability kernel under maintenance & T1 deadline contrast; T11 funded/underfunded contrast & strict viability shrinkage under collapse budget & shrinkage is proved generically, not universal empty collapse \\
\bottomrule
\end{tabular}
\caption{\textbf{Tbl-2: the C-package, compactly.} Each clause with its
standard mathematical object, its strongest landed component evidence, the
simplified certificate's proxy for it, and the inference it blocks.
Scope/price: component witnesses live on separately declared carriers; no
integrated instance exists (Section~\ref{sec:t17}).}
\label{tbl:cpackage}
\end{table}

\begin{figure}[t]
\centering
\fbox{\parbox{0.92\linewidth}{\small\emph{[Figure F2: C-package dependency
graph; rendered per the artifact figure plan.]} Nodes C1--C8 with dependency
edges; derived nodes D1--D4 attached to their source clauses; A1/A2 drawn
outside the structural boundary; the simplified certificate $\Certsimp$ shown
as a distinct inset, with no ordering arrow between it and the full package.}}
\caption{\textbf{F2: the C-package as a dependency graph.} Takeaway: the
clauses are one contract with declared interactions, not eight independent
slogans. Scope restriction: the graph shows the design target; it does not
assert an integrated instance, and no equivalence or ordering between the full
package and the simplified certificate is implied.}
\label{fig:cpackage}
\end{figure}

\subsection{Dependencies and independent axes}

The clauses are not symmetric. A forcing cascade organizes them: under
deadline and live-signal conditions, C3's self-coordinate is \emph{forced}
(Section~\ref{sec:t1}); C4 hiddenness, C6 temporal closure, and C2 strictness
are \emph{generic or declared}, holding under broad, explicitly stated
interface conditions; and C7 binding and C5 drive are \emph{contingent axes}
that a package may simply lack. The two contingent axes generate a four-cell
classifier lattice of degraded packages (T1b): binding and drive both absent
(a controller), binding without drive, drive without binding, and the full
subject cell. The cells are axis-independently witnessed classifications
(T7 witnesses each binding value and T5 each drive value on their own
carriers), not four jointly built systems; the lattice is a classificatory
schema, not a claim that these are the only possible structures.

The same synthesis packet lands a scoped recognition audit \gradeR{}: for this
catalog, every C1--C8 clause maps to composites of the declared six roles
without introducing a seventh irreducible role. The audit is scoped to this
package and the current role vocabulary; it does not prove any universal
completeness of the six roles.

\subsection{Why single ingredients are insufficient}
\label{sec:insufficient}

Each familiar single-ingredient proposal reappears inside the contract as one
clause, and each is separated from sufficiency by an explicit control
elsewhere in the catalog. Agency or regulation alone is C5/C8 material:
the credit audit exists precisely because apparent drive can be externally
scheduled or internalized without being creditable to self
(Section~\ref{sec:selfaudit}), and a nonempty viability kernel is controlled
persistence, not experience. Persistence alone (C1/C8) is exhibited by any
maintained fixed point; the T4 study shows a certificate-free reading of the
same public face (Section~\ref{sec:interface}). Integration or binding alone
(C7) collapses under marking while the branch family survives
(Section~\ref{sec:structure}). Report alone (D4) licenses only descended
boundary content, and the report map is provably non-injective on hidden
states (Section~\ref{sec:interface}). The contract is the conjunction plus its
coherence conditions; the catalog's controls are what make each ``alone''
claim fail exactly rather than rhetorically.

\subsection{Where bridges enter}

Three places, and only three, connect the structural contract to
non-structural readings. First, recognition bridges \gradeR{}: the named
proposals that D1's dual prices have the profile of felt valence
(Section~\ref{sec:structure}), that clause patterns correspond to
phenomenological categories, and that D3's marking dynamics point at
attention. Each is retractable without touching Tier~S. Second, the
ontological postulates A1/A2, spent only in Section~\ref{sec:t17}. Third, the
prospective empirical interface, which is specified as unmet obligations in
Section~\ref{sec:empirical}. Everything else in the paper is lane S.

\subsection{The full target versus the simplified certificate}
\label{sec:fullvssimp}

One distinction must be visible before any result is stated, because the
realization section depends on it. The full C-package $\Cpkg$ of
Definition~\ref{def:cpackage} is the design target. What is currently
\emph{mechanized} is a simplified, carrier-indexed certificate schema
$\Certsimp(X)$ (Section~\ref{sec:t17}): mechanics-based clause predicates
C1--C8, cross-clause coherence, and a simplified promotion ledger. The schemas
are not ordered by any landed theorem: the simplified certificate uses a
weaker imported gate ledger while adding one-step coherence restrictions
absent from the full design definition; no implication or equivalence between
them has been proved or is claimed. Three specific weakenings recur and are
flagged where used: the simplified C3 is one-step stable rather than
recursively terminal in full generality; the simplified C6 is a
budget-disciplined trace proxy rather than the full imported clock theorem;
and the simplified C8 proves strict viability \emph{shrinkage} under a
dominated collapse budget, not universal empty collapse. Table~\ref{tbl:cpackage}
carries these boundaries clause by clause.

 \section{Exact forcing of a predictive self-coordinate}
\label{sec:t1}


The catalog's strongest result answers question~1 of
Section~\ref{sec:questions}: there is a fully declared finite carrier on which
self-description through the public interface provably cannot keep the system
alive, and the unique cheapest repair is a predictive self-coordinate. This
section states the result, gives the proof idea and the decisive enumerations,
and fixes its boundary. Appendix~\ref{app:t1} contains the complete carrier,
matrices, fixed-point computations, and Lean mapping. The repository name of
this law is T1, ``the forcing of the I''; its content is an exact
predictive-self \emph{refinement} result on a declared carrier, not a
metaphysical theorem about selves.

\subsection{The deadline carrier}
\label{sec:t1carrier}

\begin{construction}[T1-L1, the self-wheel]
\label{con:t1}
The carrier is $Z = P \times W$ with $|P|=6$ public positions and $|W|=2$
hidden phases, so $N = |Z| = 12$. The public self-lens is $q(p,w)=p$: the
system can observe its own position but not its phase. The dynamics contain a
live, unobserved \emph{phase slip} with exact branch weights $3/4$ and $1/4$,
placed so that public history does not reconstruct the phase: the public trace
$(1,2,1)$ can end in either phase. A \emph{deadline} makes the hidden phase
viability-relevant: continuation beyond the deadline requires acting on the
phase-dependent continuation law, and a wrong commitment exits the safe set.
\end{construction}

Three quotients of $Z$ organize everything.
\begin{itemize}
\item The \emph{public quotient} $\Pi_q$ has the six position classes.
\item The \emph{continuation-law quotient} $L$ groups states by
viability-relevant continuation-law equality; it has five classes. $L$ is a
policy abstraction, and it is \emph{incomparable} with $q$: its classes cross
public fibers.
\item The \emph{canonical predictive quotient} $M = q \times L$, the join of
public class and continuation law, has nine classes, strictly refines $q$, and
splits exactly the public fibers at positions $1$, $2$, and $5$. Its
comparison map $M \to \Pi_q$ exists by construction.
\end{itemize}

A counting fact from the corpus's first foundations paper frames why hidden
volume makes such refinements generic: a lens with $K$ values defines exactly
$2^K$ of the $2^N$ predicates on the carrier, so here only $64$ of $4096$
predicates are $q$-definable, and a uniformly sampled predicate is
non-definable with probability $1 - 2^{-6}$. Genericity is a proxy that makes
strict extension cheap and typical when hidden states exist; the theorem below
is about the specific, viability-relevant obstruction, which the carrier
certifies by the split pair $(p_1 w_0, p_1 w_1)$: same public class, different
canonical $M$-class, hence $\Dfact(q, M) \neq \emptyset$.

\subsection{The theorem}

\begin{theorem}[Exact forcing and minimal restoration; T1]
\label{thm:t1}
\gradeS{} \emph{Grade:} theorem-grade on the exact T1-L1 carrier; conditional
family schema beyond it. \emph{Scope:} the declared twelve-state deadline
carrier, its action family, and its lens. On this carrier:
\begin{enumerate}
\item \textbf{Public exhaustion.} Every $q$-measurable feedback policy has an
empty deadline viability kernel: $\Kviab(q) = \emptyset$, by exhaustive
enumeration of all public-lens policies.
\item \textbf{Predictive restoration.} The five-class continuation-law
abstraction $L$ restores a ten-state kernel, and the canonical predictive
quotient $M = q \times L$ restores exactly the same ten states:
$\Kviab(L) = \Kviab(M)$, $|\Kviab(M)| = 10$, certified by a separate
nine-class search and a refinement audit.
\item \textbf{Unique minimal repair.} Among all $64$ refinements of the
public lens, the unique minimal closure-restoring refinement is the
seven-class quotient $\Pi_*$ that splits exactly the public fiber $p=1$;
every closure-restoring refinement splits that fiber, and the maximal
refinement that leaves $p=1$ unsplit has empty kernel.
\end{enumerate}
\emph{Assumptions:} deadline active; live, resettable, non-publicly-definable
phase. \emph{Nonclaim:} this does not show that all bounded self-predicting
systems must form selves, and it is not an empirical result about any actual
system.
\end{theorem}

Here $p=1$ is a public \emph{position coordinate} value naming the fiber the
minimal refinement splits; it is not a probability, and not a refinement
parameter.

Three distinct objects now exist, and the paper never identifies them: the
canonical nine-class $M$; the unique minimal seven-class full-carrier
refinement $\Pi_*$; and (below) the seven-class terminal quotient on the
ten-state kernel. The last two share a class count and are different
partitions of different domains.

\subsection{Proof idea and decisive enumerations}

The proof is an audit trail, not an estimate; every step is an exact
enumeration.

\emph{Exhaustion.} Enumerate every $q$-measurable feedback policy and run the
greatest-fixed-point viability iteration for each. All public-lens policies
yield the empty kernel. The obstruction is exactly the split pair: the two
histories with public trace $(1,2,1)$ end in different phases with different
continuation laws, and by the deadline the difference is lethal to one
commitment.

\emph{Restoration.} The five-class $L$ implements the policy
(need-clockwise $\to$ clockwise, need-counterclockwise $\to$
counterclockwise, return $\to$ stay) and restores a ten-state kernel; the
nine-class $M$ search returns the same ten states. Note the logical roles:
$L$ is the policy abstraction that suffices, $M$ is the canonical object that
carries it while remaining comparable with the public lens.

\emph{Minimality.} Audit the six single-fiber splits of $q$: split kernel
sizes are $0, 10, 0, 0, 0, 0$ across the fibers $p = 0, \dots, 5$, so only the
$p=1$ split restores the kernel. The maximal refinement leaving $p=1$ unsplit
has empty kernel; since a finer lens can implement every coarser-lens policy,
the best kernel never shrinks under refinement, so every refinement that
leaves $p=1$ unsplit is empty. Uniqueness and minimality of the seven-class
$\Pi_*$ follow over the full $64$-refinement lattice. (The monotonicity step
is a prose lemma over the eight computed searches; the enumeration itself is
Python-certified, as flagged below.)

\emph{Controls.} Two controls show the conditions are load-bearing. Removing
the deadline makes $\Kviab(q)$ nonempty with the public policy \emph{stay}:
no forcing without stakes. Removing the live slip (the washout control) makes
the phase publicly reconstructible from the calibrated start, and the audit
then classifies the hidden coordinate as an artifact rather than a forced
self-coordinate: a merely decodable correlate is refused.

\subsection{Closure objecthood and the terminal self-lens}
\label{sec:t1c1c3}

Two further exact computations land the C1 and C3 clauses on the same
carrier's self-package, the ten-state kernel under the canonical $M$
policy.

\emph{C1 (objecthood).} With $P$ the closed-loop Markov matrix, $C$ the
$M$-class indicator, and $U$ the uniform prototype on realized $M$-fibers,
the staged closure is $E_\tau = P^\tau C U$. Exact Gaussian elimination gives
the horizon table: at $\tau = 1,2,3,4$ the idempotence defects are $1$,
$1/8$, $1$, $3/32$, with unique CORE-uniform fixed points at odd horizons and
a two-extremal fixed segment at even horizons (the segment's midpoint is the
odd-horizon fixed point). The odd-horizon defect $1$ is recorded plainly:
no smallness is claimed, and the fixed distributions are declared proxy
objects.
Robustness is a separate audit at $\tau=2$: mixing $P$ with the uniform
kernel at $\varepsilon = 1/16$ collapses the segment to a \emph{unique}
nonuniform fixed point that retains CORE mass $77/80 \ge 3/4$, sits at total
variation $3/80$ from CORE-uniform, and has exact defect $78795/524288$; a
second row at $\varepsilon = 1/100$ gives total variation $3/500$, evidencing
$\mathrm{TV} = (3/5)\varepsilon$ at both tested points.

\emph{C3 (terminal self-lens).} Start from the public partition restricted to
the ten-state kernel (six classes) and repeatedly split every class by its
exact per-action transition profile, with out-of-kernel exits labeled
distinctly. The size trajectory is $6 \to 7 \to 7$: the public partition is
unstable at the pair $(p_1 w_0, p_1 w_1)$, one refinement pass reaches
stability, and all $100$ ordered-pair checks identify the terminal partition
with canonical $M$ restricted to the kernel. Extensionally
$M(z) = (q(z), L(z))$: the terminal self-lens carries the state's own packaged
continuation coordinate. That is the exact, finite sense in which the carrier
is forced to know where it is \emph{in its own continuation law}, not merely
where it is publicly.

\begin{table}[t]
\centering\small
\begin{tabular}{lr}
\toprule
Quantity & Exact value \\
\midrule
Carrier size $|Z|$; public values & $12$; $6$ \\
Predicates definable from $q$ / total & $64$ / $4096$ \\
Public deadline kernel $|\Kviab(q)|$ (all policies, exhaustive) & $0$ \\
Continuation-law kernel $|\Kviab(L)|$ & $10$ \\
Canonical predictive kernel $|\Kviab(M)|$ (${}=\Kviab(L)$, audited) & $10$ \\
Refinements of $q$ exhaustively checked & $64$ \\
Unique minimal restoring refinement $\Pi_*$ (splits fiber $p=1$) & $7$ classes \\
Public partition on the ten-state kernel & $6$ classes \\
Terminal refinement on the kernel (${}=M$ restricted, $100$ pair checks) & $7$ classes \\
Closure defects $\delta_{\tau}$, $\tau=1,\dots,4$ & $1,\ 1/8,\ 1,\ 3/32$ \\
Perturbation $\varepsilon$; CORE mass retained & $1/16$; $77/80$ \\
Total-variation witness at $\varepsilon=1/16$ & $3/80$ \\
Aggregate defect at $\varepsilon=1/16$ & $78795/524288$ \\
No-deadline control & $\Kviab(q)\neq\emptyset$ \\
Washout (no-slip) control & phase reconstructible; not forcing \\
\bottomrule
\end{tabular}
\caption{\textbf{Tbl-3: the T1 exact result card.} Computed rows are exact
rationals from the frozen result snapshot; the two counting rows (predicates
$2^6$/$2^{12}$; the $64$-refinement audit scope) are definitional counts
stated by the theorem catalog. Scope/price: the declared
twelve-state carrier only; the full-carrier minimal seven-class split and the
kernel-restricted seven-class terminal quotient are distinct objects despite
the shared class count.}
\label{tbl:t1}
\end{table}

\begin{figure}[t]
\centering
\fbox{\parbox{0.92\linewidth}{\small\emph{[Figure F3: T1 forcing and
refinement; rendered per the artifact figure plan.]} Three aligned panels:
(a)~the twelve-state carrier with actions, public positions, and the hidden
phase slip; (b)~public versus predictive/refined viability kernels ($0$ versus
$10$ states); (c)~the refinement lattice with the unique minimal seven-class
split of the $p=1$ fiber, drawn separately from the kernel-restricted
trajectory $6 \to 7 \to 7$.}}
\caption{\textbf{F3: exact forcing on the declared carrier.} Takeaway: public
self-description fails exhaustively; one minimal predictive split restores
viability. Scope restriction: exact declared carrier; the family extension is
conditional and is not plotted as established.}
\label{fig:t1}
\end{figure}

\subsection{Exact carrier versus family schema}
\label{sec:t1family}

What generalizes, and at what price, is stated once and inherited everywhere.
For substrates satisfying the published deadline-family conditions (a live,
resettable, non-publicly-definable phase, and a deadline that makes the hidden
continuation law viability-relevant), self-prediction through a bounded
public self-lens generically carries a predictive non-factorization
obstruction, and under the deadline that obstruction is viability-relevant.
That family statement is a \emph{conditional schema} \gradeS{}: all
hypotheses displayed, none discharged beyond the exact carrier. The
conditions are strong, and systems lacking them fall into the controls, not
into the forced case. No claim is made that the conditions are universal or
empirically common, and no generic ``selfhood theorem'' is asserted.

\emph{Mechanization boundary.} The Lean kernel for this law formalizes the
full twelve-state carrier and the finite public, continuation-law, and
canonical quotient tables; the exhaustive policy and refinement searches are
certified by the Python study and its tests, not by a general Lean fixed-point
theorem. Lean's C1 facts cover the $\tau=1$ operator (fixed-point verification
by evaluation and the defect witness), not uniqueness, the $\tau=2..4$ rows,
or the perturbation audit; Lean's C3 facts prove terminal-partition stability
and public instability with the $6/7/7$ sizes as declared constants, while the
refinement execution is Python-certified.

\subsection{Position among precursors}
\label{sec:t1position}

The ingredients are standard and are credited as such; the composition and
the forcing question are not found in the precursors.

\emph{Computational mechanics.} Causal states quotient semi-infinite
stationary histories by conditional future distributions, with sufficiency,
minimality, and uniqueness theorems under those hypotheses
\citep{CrutchfieldYoung1989, ShaliziCrutchfield2001}. T1's quotients are the
finite controlled analog, but the problem is different: a finite deadline
carrier with feedback control and a viability criterion, not a stationary
process. In particular, causal-state minimality cannot prove
Theorem~\ref{thm:t1}(3): the minimal object there is a \emph{refinement of a
declared public lens} restoring a control property, not a minimal sufficient
statistic. That is why the canonical $M$ has nine classes while the minimal
restoring refinement has seven.

\emph{Viability theory.} The kernel, its greatest-fixed-point
characterization, and the descending iteration are the standard discrete-time
constructions \citep{AnayaEtAl2009Viability, AubinBayenSaintPierre2011}. The
paper's addition is what is being iterated over: lens-measurable policy
classes, so that kernels become properties of \emph{interfaces}.

\emph{Regulation.} The good-regulator tradition
\citep{ConantAshby1970} shows, under its entropy-minimizing setup, that a
simplest optimal regulator implements a mapping from reguland states to
actions. That is a much narrower state-to-action result: it does not
establish an internal self-model, a predictive quotient, a viability kernel,
or an exact public-refinement minimality, and nothing here treats it as if it
did.

 \section{Structural consequence clusters}
\label{sec:clusters}


The remaining laws are organized consequences and exact witnesses, not
eighteen equal headline novelties. Each subsection gives the setup shape, the
central result with its grade, the proof idea needed for the conceptual
consequence, and the controls; Table~\ref{tbl:catalog} collects the catalog
compactly, and Appendix~\ref{app:compendium} holds every remaining exact
value. Throughout, each law's repository name (``privacy of qualia,''
``zombie,'' ``free will,'' \dots) labels the classical problem it recasts; the
restricted content follows the name immediately, every time.

\subsection{Interface and access}
\label{sec:interface}


Four laws share one diagram (Figure~\ref{fig:interface}): a hidden object, a
public/report quotient, and an attempted factorization whose exact success or
failure, always relative to a declared interface, is the content.

\paragraph{T2: hiddenness at a declared interface.}
Fix a declared public channel family $\Lambda$, the current public quotient
$Q$ it induces, the determining predictive quotient $M$, and the comparison
map $\pi : M \to Q$. \emph{Exposure} holds when some admissible current
observable factors through the determining state non-constantly on the
realized support. The imported hiddenness law (Appendix~\ref{app:imports})
then gives a dichotomy controlled entirely by the binary exposure predicate,
not by coupling strength.

\begin{theorem}[Interface-relative non-exposure; T2]
\gradeS{} \emph{Grade:} theorem-grade direct finite instantiation; the C4
condition is generic and declared, not universal. \emph{Scope:} the declared
channel family $\Lambda$ on the finite privacy-wheel carrier. Under
non-exposure with predictive witness content, there is a connector loop that
fixes every current public class while moving the predictive class. On the
exact carrier: two hidden states, one public class; the loop moves $0/1$ of
$Q$-classes and $1/1$ of $M$-classes. Adding a public $M$-reading channel
dissolves the witness (the dissolution control); the degenerate $M=Q$ control
has no witness and licenses no privacy claim. \emph{Nonclaim:} nothing is
private under every possible interface, and hiddenness does not imply
phenomenality.
\end{theorem}

The classical resonance is deliberate: ``what it is like'' talk fixes an
explanandum of subjective character and point of view \citep{Nagel1974}.
T2's contribution to that explanandum is exact and modest: it turns
``private'' into a computable, channel-relative predicate with a dissolution
control. Add one instrument, and the same carrier's privacy verdict
flips.

\paragraph{T3: the explanatory gap as a strictness certificate.}
Strict promotion means a nonempty factorization defect. The imported
nonfactorization dichotomy says $\Dfact(\pi_0, \pi_1)\neq\emptyset$ if and
only if $\pi_1$ does not factor through $\pi_0$.

\begin{theorem}[No-shortcut corollary; T3]
\gradeS{} \emph{Grade:} corollary-grade. \emph{Scope:} declared source and
target maps on a fixed finite interface. If the promoted content is strict
over the public layer, no lower-layer shortcut derivation recovers it: the
factor map does not exist, and any closed lower-layer compression fails
closure for the promoted readout. Contrapositive: if a proposed derivation
succeeds, the extension was never strict. On the exact carrier (reusing the
T9 support below), the strict case computes factor attempt $=$ none from the
shared-report pair; the control with report as both source and target
computes a concrete factor map and no gap. \emph{Nonclaim:} this neither
solves nor dissolves the philosophical explanatory gap in general, asserts no
hidden substance, and does not say mechanism is absent: if a derivation
succeeds, the content was non-strict, not unmechanistic.
\end{theorem}

The reading this licenses, and the only one, is structural: the
explanatory gap, in the finite formal sense, is the \emph{audit certificate of
strictness}, not a refutation of anything. The historical label is credited to
\citet{Levine1983}; the hard/easy framing to \citet{Chalmers1995FacingUp}.
Whether the philosophical gap has this structure is an interpretive question
the theorem does not settle.

\paragraph{T9: reports license descended content only.}
Model experience as a hidden promoted object $U$ that the declared lower
interface cannot accept directly; a report channel is a boundary-expression
former with an audited descent. The imported no-overreading theorem says an
accepted claim has an empty overreading defect: every suppressed use is
admissibly bridged.

\begin{theorem}[Report descent and no-overreading; T9]
\gradeS{} \emph{Grade:} direct instantiation of imported audit theorems plus
finite toy witness. \emph{Scope:} declared report and no-report interfaces on
a three-state hidden support. A valid report licenses only its descended
boundary object, never the hidden structure. The finite content: the report
map is non-injective (states $u_0 \neq u_1$ share the report boundary
$r_{\mathrm{shared}}$), while a differently scoped no-report instrument
distinguishes the same pair; an overreaching claim (suppressed unbridged use)
classifies as blocked, and a visible descent-only claim classifies as
accepted, under the same vendored classifier. \emph{Nonclaim:} reports are
not useless or disconnected; and hidden content is not inaccessible to all
possible instruments, since the no-report instrument in the same study
distinguishes what the report identifies.
\end{theorem}

In this vocabulary, the access/phenomenal distinction \citep{Block1995} is
layer-structural: access is D4, the descended report channel; phenomenal
structure is the hidden $U$. The mapping is comparative, not an identity
claim, and no measurement of phenomenal consciousness is asserted. Every
report-level claim elsewhere in this paper, including every future
empirical-interface row, inherits T9's boundary: descent only, unless a
further bridge is separately audited.

\paragraph{T4: zombies as certificate contrast on one public face.}
The zombie study has three separated parts.

\begin{theorem}[Leibniz quotient; T4a]
\gradeS{} \emph{Grade:} theorem-grade direct quotient instantiation.
\emph{Scope:} the declared admissible public lens family at one layer. Any
legitimate layer-level quantity factors through the quotient by empirical
indiscernibility; duplicates indiscernible under \emph{all} admissible lenses
are one object of that layer. ``Zombie'' is therefore not a public-layer
distinction when all declared public lenses agree. \emph{Nonclaim:} sameness
is layer-relative; there is no metaphysical identity claim at every layer, and
empirical equivalence does not prove phenomenal identity.
\end{theorem}

\begin{theorem}[Certificate demarcation; T4b]
\gradeS{} \emph{Grade:} schema-grade certificate contrast. \emph{Scope:} one
public carrier face with an imported certificate battery. A subject/zombie
distinction exists relative to a declared audit only if the certificate lens
family separates the duplicates. The exact study realizes this as two
\emph{certificate readings of one} public carrier face: reading A carries the
genuine imported T2/T7/T5/T6 witnesses (hiddenness, strict recombination,
self-credited drive, same-level-audit content); reading B carries the
corresponding degenerate, collapsed, and no-drive controls. Public statistics
match exhaustively over the declared current observable; an imperfect-match
control C, which adds a public $M$-reading channel, is classified publicly
discernible \emph{before} any certificate comparison. \emph{Nonclaim:} A and
B are not two independently constructed integrated carriers, and no
metaphysical zombie detector is offered.
\end{theorem}

T4c, the reading that achieving the closures is realization rather than
imitation, is \gradeO{} and is deferred entirely to the conditional bridge
of Section~\ref{sec:t17}; the interface cluster uses only the one-sentence
pointer.

\begin{figure}[t]
\centering
\fbox{\parbox{0.92\linewidth}{\small\emph{[Figure F4: interface-relative
opacity; rendered per the artifact figure plan.]} One commuting/noncommuting
factorization diagram shared by T2, T3, T4, and T9: hidden object $U$ /
determining quotient $M$; public and report quotients; the attempted factor
map (crossed out where the defect is nonempty); the T4 public face with its
two certificate readings; and the exposure/dissolution control as an added
channel.}}
\caption{\textbf{F4: one diagram, four laws.} Takeaway: privacy, the gap,
report limits, and zombie contrasts are the same factorization question asked
at four interfaces. Scope restriction: every failed factorization carries its
declared interface; no absolute privacy and no integrated zombie pair.}
\label{fig:interface}
\end{figure}

\subsection{Self-audit, action credit, and the meta-problem}
\label{sec:selfaudit}


\paragraph{T5: a priced free-will construction.}
The free-will law is four structural clauses and two foreclosures on one
deterministic deliberation carrier; its repository name promises no
metaphysical freedom theorem, and none is delivered.

\begin{theorem}[Priced residue and audited authorship; T5]
\label{thm:t5}
\gradeS{} \emph{Grade:} clauses graded individually below. \emph{Scope:} one
declared deterministic support, bounded observer budgets, declared generator
provenance. \emph{Nonclaim:} no ontic indeterminacy, no metaphysical
freedom, no disproof of every libertarian view.
\begin{enumerate}
\item \emph{Lawfulness} (framing). Choices are readouts of the run; no
spontaneity primitive is added.
\item \emph{External underivability} (theorem-grade toy instantiation). A
bounded public observer at budget $B$ can face positive priced residue:
same-summary states with different deterministic choices. The exact scores at
budgets $0,1,2$ are $1, 1, 1/2$, non-increasing, with a split witness at every
budget.
\item \emph{Self-underivability} (direct imported-theorem instantiation). The
self-model is a bounded packaging; a complete same-level self-audit
classifies \texttt{undefined\_circular} (an out-of-scope variant:
\texttt{outside\_scope}), and the self-model summary has its own
same-summary/different-choice split pair.
\item \emph{Authorship} (direct instantiation of no-fake-arrow audits).
Directionality is credited to the agent only when it survives the
protocol-trap and data-processing nulls and one common credit rule returns
\emph{self}: genuine autonomous drive has unequal exact self-channel cycle
products $8/27$ versus $1/27$; the externally scheduled apparent arrow has
unequal \emph{other}-channel products ($0$ versus $16/81$), a real arrow
credited to other rather than to self; and after schedule internalization the lifted kernel decomposes
exactly into clock and content channels, both balanced ($0$ versus $0$ on
credited products; canonical lifted cycle $1/324$ each way), so the same rule
returns \emph{none}. Deterministic coarsening of the blocked case leaves
equal path fluxes ($1/27$ each way): observation does not create affinity.
\item \emph{Libertarianism foreclosed} (foreclosure argument). The positive
residue of clauses 2--3 is bounded-observer ignorance under declared
packaging and budget, not substrate spontaneity.
\item \emph{Deflation foreclosed} \gradeO-flavored. ``Only a lower-layer
event'' commits the factorization error: once the layer is strict, the layer
claim cannot be replaced by a base-layer paraphrase without losing the
certificate, as T3 establishes. This is a foreclosure argument with an
ontological flavor, not a Tier~S theorem.
\end{enumerate}
\end{theorem}

The philosophical shape is symmetric by design: the same exact carrier
prices both the libertarian reading (residue is ignorance, not spontaneity)
and the hard-determinist deflation (paraphrase loses a certified
distinction). What remains is an audited, model-relative credit assignment
(self, other, or none) on declared channels, and nothing stronger.

\paragraph{T6: introspection limits without introspection nihilism.}
\begin{theorem}[Same-level audit failure, rotating towers, confabulation; T6]
\gradeS{} \emph{Grade:} direct instantiation of imported self-audit theorems
plus finite toy extension. \emph{Scope:} instrument-indexed claims; complete
same-level self-certification, not ordinary introspection. No instrument
certifies its own complete soundness at its own level: the claim classifies
\texttt{undefined\_circular} or \texttt{outside\_scope}; a level-shift bridge
reclassifies it as \texttt{provisional}, never as same-level final. Every
finite rotating-audit tower has an uncertified top: the three-level toy tower
computes own-level statuses (\texttt{undefined\_circular} at each level),
bridge statuses (\texttt{provisional}, \texttt{provisional}), and
top-requires-next-level $=$ true. \emph{Confabulation corollary} (toy
extension; schema beyond it): with an audit update missing, the unique
closure-preserving completion from visible records is \emph{left} while the
hidden true state is \emph{right}: coherent and wrong. With the update
present, the unique completion is correct and the verdict disappears.
\emph{Nonclaim:} introspection is neither impossible nor worthless; partial,
level-shifted, bounded audit remains admissible, and self-image stability
over time is named as dependent on an open conjecture, not claimed.
\end{theorem}

The slogan form, that the I doing the auditing is never the I being audited,
is exactly as strong as the classifier statuses above, and no stronger.

\paragraph{T16: the meta-problem as a convention-relative composition.}
The meta-problem asks why systems like us judge there to be a hard problem
\citep{Chalmers2018MetaProblem}. The composition law computes one exact,
deliberately small version.

\begin{theorem}[Meta-problem composition; T16]
\gradeS{} \emph{Grade:} composition-grade and convention-relative.
\emph{Scope:} one two-state composite carrier whose public interface and
report channel coincide; an explicitly declared completion convention. On the
composite package: the T2-shape privacy witness and T3-shape strict-factor
absence hold reflexively of the package's own self-model (hidden moved
fraction $1$, public moved fraction $0$, no factor $q \to m$); the T6-shape
same-level certification of its report channel is blocked
(\texttt{undefined\_circular}); and the T9-shape classification licenses the
report as a valid descent while blocking overreading. The declared completion
convention maps the computed flag triple (gap witness, audit blocked, factor
absent) $=(\top,\top,\top)$ uniquely to the report
\texttt{cannot\_derive\_own\_experience}; the exposed-interface control
computes $(\bot,\top,\bot)$ and uniquely selects
\texttt{derivation\_certified}. \emph{Nonclaim:} the generic claim that
bounded reflexive subjects report a hard problem is \emph{not} derived from
the four component laws alone; it holds relative to the declared
convention. Moreover, on this carrier the public and report coordinates coincide,
so several nominally distinct premises extensionally reduce to one
nonfactorization fact (premise independence would need a carrier that
separates them).
\end{theorem}

Under that same convention, the composition threads the intended contrast
without deleting either side's evidence: illusionism is right that reports
under-determine the hidden object, wrong if it deletes the object;
mysterianism is right that no shortcut derivation exists, wrong if it treats
the gap as failure rather than as the certificate of strictness. T16 is one
convention-relative finite composition study; it does not settle the
meta-problem.

\subsection{The structure of experience}
\label{sec:structure}


Five laws and a temporal cluster give exact finite structure to binding,
temporal experience, valence, quality geometry, and inversion, always as
declared-carrier structure and never as phenomenal measurement.

\paragraph{T7: binding as strict recombination refinement.}
Over a declared branch family (modality/content branches), two quotients
compete: the \emph{branchwise} quotient $\Kbw$, the strongest state obtainable
by treating branches independently, and the \emph{recombination} quotient
$\Rrec$, induced by admissible joint observables on assemblages, with a
forgetful map $\eta : \Rrec \to \Kbw$.

\begin{theorem}[Binding witness with marking and erasure; T7]
\gradeS{} \emph{Grade:} theorem-grade finite instantiation of imported
strict-refinement machinery. \emph{Scope:} the declared branch family and
interfaces. Binding is certified by strict refinement $\Rrec \succ \Kbw$: two
assemblages in one branchwise class and distinct recombination classes. On
the exact carrier: the coupled regime has $\Kbw$ blocks $(2)$ and $\Rrec$
blocks $(1,1)$ with witness (\emph{bound-diagonal}, \emph{bound-crossed});
\emph{marking} a branch-identifying observable collapses the partitions to
exact equality ($\Rrec = \Kbw$, both $(1,1)$); \emph{erasure} of the mark
restores the original witness; and the aggregation control, in which no joint
observable is ever declared, has $\Rrec = \Kbw$ from the start (a
never-coupled product, not a marked collapse). \emph{Nonclaim:} no derivation
of phenomenal unity or of what unity feels like, and no attention cost:
the law contains no budget, price, or scheduling variable, so it establishes
marking-collapse, not a priced attention mechanism.
\end{theorem}

\paragraph{T8: the specious present as a funded triple.}
Temporal structure is treated as three separable properties of a closure:
order (a successor), tick (a staged carrier that persists \emph{and}
progresses), and arrow (a monotone accounting variable). Each must
be maintained under budget.

\begin{theorem}[Purchased clock and no unified now; T8]
\gradeS{} \emph{Grade:} theorem-grade finite instantiation of imported clock
and protocol criteria. \emph{Scope:} the declared three-phase clock carrier
and protocol family. The experienced-present proxy is the package's
stabilized (order, tick, arrow) triple, with width the closure timescale
$\tau$. The single mechanism $\mathrm{advance} = b/(b+1)$ at budget $b=2$
gives exact progress $2/3$ and tick rate $2/9$; the funded nine-step trace
has six advances, two completed wraps, and a monotone ledger with
$\Delta A = 2 > 0$; the unfunded control ($b=0$) stalls exactly ($0$
progress, $0$ ticks, constant trace, no arrow). No single global ``now'' is
guaranteed across sub-closures: the noncommuting protocol pair (rotation,
swap) has holonomy defect $1$ and the commuting control ($R$, $R^2$) defect
$0$, so path-dependent local time translations admit no global potential.
Finally, felt-flow candidates route through the audit channel, never through
holonomy alone: the holonomy-without-drive control derives an unbiased audit
support with equal products ($3/1000$ each way) and is classified blocked.
\emph{Nonclaim:} no universal specious-present width, no physical time or
spacetime, and holonomy is never a directionality certificate.
\end{theorem}

\paragraph{T11: constraint prices, and valence as a named bridge.}
Any package with C8 maintenance constraints is budgeted; closure under
binding constraints yields exact dual prices $\lambda_i = \partial \Phi^* /
\partial b_i$.

\begin{theorem}[Exact dual prices; T11 core]
\gradeS{} \emph{Grade:} theorem-grade finite duality core. \emph{Scope:} one
exact rational linear program over a three-action policy simplex with two
maintenance budgets. At binding budgets $b=(2,2)$: optimum value $76/11$,
both constraints tight, exact local duals $\lambda = (24/11,\, 14/11)$. Slack
regime at $b=(5,1)$: $\lambda_1 = 0$ exactly. Tightening the first budget
across $(4,3,2,3/2)$ recomputes $\lambda_1 = (0,\, 4/3,\, 24/11,\, 24/11)$,
non-decreasing. Proxy discipline: a declared row-permuted proxy cost yields
unstable recovered proxy prices across two environment phases ($4$ and $0$),
while the true LP is phase-independent by construction; a planted dual
certificate with positive price on a slack constraint violates complementary
slackness and is rejected. Viability tie: the kernel is nonempty when both
budgets meet their thresholds and empty when either is underfunded by
$1/11$. \emph{Nonclaim:} shadow-price duality is standard mathematics, and
none of these values generalizes to actual subjects.
\end{theorem}

\begin{remark}[The valence bridge]
\gradeR{} The identification of the dual-price profile with \emph{felt
valence} is a recognition bridge, named here and nowhere derived: the C8 dual
has the intended profile of graded magnitude, sign/pressure, motivational
relevance, and a direct tie to viability. A counterexample to the identification
retracts the bridge, not the duality mathematics. Felt valence is never
Tier~S in this paper.
\end{remark}

\paragraph{T12: a descended quality geometry, exactly audited.}
\begin{theorem}[Induced-kernel similarity geometry; T12]
\gradeS{} \emph{Grade:} theorem-grade induced-kernel construction and exact
audit; the geometry-birth grade is audit-only. \emph{Scope:} the corrected
32-state micro-carrier (eight hues, two chroma, two texture levels), packaging
lens forgetting texture, staged twice. The lens descends the symmetric micro
kernel to a 16-point macro kernel $\hat P = U P^2 C$ with exact lift
consistency $UC = I$; texture is deliberately macro-relevant (a sluggish
grain multiplies move rates by $1/2$), so prototype content is operational,
not decorative. The exact max-product path similarity gives warm-neighbor
similarity $43/270$ versus cool-target $1/810$, so the warm neighbor is
strictly nearer; $d = -\log D$ is the additive distance on positive support.
Audits, all exact: closure defect $13441/52488$ (the swapped-prototype
incoherent control: $2619989/3499200$, strictly larger); prototype stability
maximum $2/3$, mean $13801/25920$; all $256$ ordered pairs connected; all
$16^3$ multiplicative triangle checks pass; route mismatch $0$. A texture-0
point-prototype lift preserves $UC=I$ yet changes the induced object and is
rejected: lift consistency alone does not certify prototype content. A local
signature table is chroma-favored at hues $0,3,7$, balanced at $1,5,6$,
hue-favored at $2,4$; a hand-declared signature-change transport has loop
residue $0$ on the balanced plateau and $1$ at the warm/cool boundary, with
multiplicative residual ratios $180/169$ (plateau) and $102/59$ (boundary),
both above one and strictly ordered. \emph{Nonclaim:} no smallness threshold
adjudicates the audits (only the computed controls do); the full
geometry-birth checklist (fitted inter-scale distortion, dimension
diagnostics, and curvature via neighborhood-MDS/Procrustes transport) is
not implemented; geometry is layer-relative; and the recovered similarity is
descended structure, never the hidden experience.
\end{theorem}

\paragraph{T13: the inverted spectrum as a gauge computation.}
\begin{theorem}[Descent-fixing gauge and rigidification; T13]
\gradeS{} \emph{Grade:} theorem-grade exact finite construction; the broader
quality-space rigidification principle is schema-grade. \emph{Scope:} an
eight-state fiber bundle, $|Q| = 4$ public classes with two fiber points
each, declared descents. Define the descent-fixing gauge group
$G = \{\sigma \in \mathrm{Aut}(M) : \pi\sigma = \pi \text{ and } d\sigma = d
\text{ for every admissible descent } d\}$. Spectrum inversion is a lawful
distinct possibility exactly to the extent that $G$ is nontrivial. Exact
enumeration: $|\mathrm{Aut}|=40320$ automorphisms of the support, $16$
projection-fixing, and $G = \{\mathrm{id}, \text{global swap}\} \cong
\mathbb{Z}_2$ in the symmetric case; adding \emph{one} declared,
descent-visible discrimination-density orientation (calibrated by T12's
landed local signature at one public class) recomputes $G = \{\mathrm{id}\}$.
Route transport between two isomorphic fibers shows the same contrast: two
admissible transports symmetric, exactly one rigidified. \emph{Nonclaim:} no
claim that any particular quality space has nontrivial $G$; the rigidifier is
a modeled assumption, not a measurement; and if richer quality spaces always
rigidify, the inverted spectrum dissolves, which is an allowed outcome.
\end{theorem}

The philosophical translation is precise: ``is your red my red'' has
determinate content exactly where the admissible descents rigidify the route;
where $G$ remains nontrivial, the residual comparison is lawfully
gauge-indeterminate rather than merely unreported.

\paragraph{T18: the temporal consequence cluster.}
Four thin, exact extensions of T8, plus one demotion, complete the temporal
story.

\begin{itemize}
\item \textbf{T18a (layer-minted proper time).} \gradeS{} theorem-grade finite
extension. The package's temporal coordinate does not factor through
substrate time: the exact strictness witness is a state pair with equal
substrate time and different package tick parity, so
$\Dfact(\text{substrate time}, \text{package time}) \neq \emptyset$; T8's
funded triple is reused unchanged. Nonclaim: no relativistic proper time, no
claim that every subject has this clock.
\item \textbf{T18b (enablement as a lens rewrite).} \gradeS{} theorem-grade
exact finite instantiation. Promotion of a layer is an \emph{enablement-time}
event: a closure rewrite that creates the variable set in which the layer's
causation-time then runs. Exactly: the phase-omitting coarse lens has closure
defects $(1,1,1)$ across horizons; the enriched lens has $(0,0,0)$; the
partial-enrichment variant drops $(4/7,4/7,4/7)$ to $(0,0,0)$ using only the
relevant parity coordinate, strictly coarser than full reveal; the no-birth
control has zero defects before and after, so no beneficial rewrite exists.
Nonclaim (the repository name is ``subject birth''): no subject is created or
detected; the theorem is a finite lens-rewrite computation.
\item \textbf{T18c (demoted).} The intended theorem, temporal closure as a
precondition for stating the other clauses, is \emph{not} landed. The
static source audit finds explicit horizon/deadline tokens only in T1 among
T1/T2/T5/T7. T18c is therefore a schema-grade dependency observation,
recorded as demoted, and nothing in this paper cites it as a theorem.
\item \textbf{T18d (state taxonomy by funding and gating).} \gradeS{}
taxonomy-grade finite construction. Funding and entrainment gates classify
temporal-state rows: unfunded (tick rate $0$); funded and entrained (match
rate $1/1$, tick $1/1$); funded but decoupled (match $1/3$, tick $1/1$). The
rows are descended structural readings under T9's report discipline:
no claim about actual sleep, dreaming, or clinical states is licensed without
a recognition bridge, and none is proposed here.
\item \textbf{T18e (no canonical now-transport).} \gradeS{} direct temporal
reading of T13's machinery. Comparing ``nows'' across sub-closures is a route
transport problem: two admissible transports in the symmetric case, exactly
one after rigidification. There is no canonical inter-layer now unless the
declared descents rigidify the transport.
\end{itemize}

\begin{figure}[t]
\centering
\fbox{\parbox{0.92\linewidth}{\small\emph{[Figure F5: structural
consequence atlas; rendered per the artifact figure plan.]} Four exact
panels: (a)~T7 marking/erasure partitions ($\Kbw$/$\Rrec$ block structures
across the four regimes); (b)~T8 funded versus unfunded tick trace with the
monotone ledger; (c)~T11 dual-price tightening curve $0, 4/3, 24/11, 24/11$
with the slack regime marked; (d)~T13 gauge and transport counts, $2 \to 1$
under rigidification. The felt-valence label appears outside panel (c),
marked Tier~R.}}
\caption{\textbf{F5: four exact structural consequences.} Takeaway: binding,
temporal funding, constraint prices, and gauge rigidification are computed
objects on declared carriers. Scope restriction: panel values are the exact
snapshot values; the T8 arrow is a stepped trace, not a rate.}
\label{fig:atlas}
\end{figure}

\subsection{Persistence and boundaries}
\label{sec:boundaries}


\paragraph{T10: a graded, hysteretic, clause-decomposable boundary.}
\begin{theorem}[Existence boundary schema; T10]
\gradeS{} \emph{Grade:} schema plus exact correspondences. \emph{Scope:} a
three-state maintenance carrier with a threaded record accumulator (credit
accrues $1/8$ after a nonempty kernel to a cap of $1/4$, decays $1/8$ after
an empty one) over the budget grid $\{0, 1/4, 1/2, 3/4, 1\}$.
Existence-boundary labels are modeled as viability collapses, that is, as
crossings of $\Kviab=\emptyset$ under the declared ledger, and the finite
prediction is threefold. \emph{Graded:} kernel sizes pass through intermediate
islands, up-sweep $[0,0,1,2,3]$, down-sweep $[3,3,2,1,0]$. \emph{Hysteretic:}
the up-sweep first-nonempty threshold is $1/2$ while the primed down-sweep
last-nonempty threshold is $1/4$; recovery and collapse thresholds differ
because the mechanism carries record state, not because thresholds were
tabled separately. \emph{Clause-decomposable:} C1-style persistence can
coexist with collapsed C7 binding (covert-awareness pattern) or collapsed C6
temporal closure (timeless pattern), by importing the respective controls;
full collapse zeroes all three. Robustness is conditional, not universal: a
weaker accumulator rule loses strict hysteresis and a stronger one never
re-crosses on the tested grid. \emph{Nonclaim:} no claim that actual loss or
recovery of consciousness is hysteretic, and any external measure of closure
depth would need its own descent audit, and none is given.
\end{theorem}

\paragraph{T14: identity as predictive lineage.}
\begin{theorem}[Lineage, trauma, fission/fusion, empty target; T14]
\gradeS{} \emph{Grade:} theorem-grade transfer of imported route-transport
machinery plus finite extensions. \emph{Scope:} a declared memory-wheel route
package with three histories, one current class, two predictive classes.
Identity over time is represented by the lineage of the predictive quotient
$M$: by the imported reachable-image factorization theorem, any
future-sufficient self-summary factors to $M$ (the toy computes the factor
from a strictly finer candidate and rejects the reverse). \emph{Trauma
shape:} a loop that is current-trivial and predictively nontrivial moves
$0/1$ of current classes and $1/1$ of predictive classes: same by all
present tests, changed in continuation. The identity-loop control moves
neither. \emph{Fission/fusion:} continuation maps on $M$-fibers classify
exactly as fiber-splitting (duplicated support) and fiber-merging (merged
support). \emph{Empty target:} the funded support has kernel size $3$; the
underfunded support has $\Kviab=\emptyset$ and empty post-exit fiber image,
so no continuation preserves the pre-exit $M$-fiber. \emph{Nonclaim:} no
complete metaphysics of personal identity. Token-independence and
death-as-dissolution readings are one-sentence Tier~O pointers deferred to
Section~\ref{sec:t17}; the Tier~S death fact is only $\Kviab=\emptyset$ plus
no fiber-preserving continuation, and the empty-target computation is
Python-certified (it is outside the law's Lean kernel).
\end{theorem}

\paragraph{T15: combination as a two-gate audit on one carrier.}
\begin{theorem}[Combination gate; T15]
\gradeS{} \emph{Grade:} schema with an exact one-carrier construction.
\emph{Scope:} one joint-package carrier with live states
\emph{diagonal}/\emph{crossed} and a dead exit. Combination into a joint
subject is certified only when the joint family passes both gates on the
\emph{joint} carrier: (i)~C7 joint strict refinement, where the exact witness has
$\Kbw = (2)$, $\Rrec = (1,1)$; and (ii)~C3 joint reflexivity, where per-action
refinement of the joint public lens terminates at the pairing-valued
predictive quotient, trajectory $1 \to 2 \to 2$, and that quotient carries the
viability-relevant joint coordinate (public-lens kernel $\{$diagonal$\}$
versus terminal-lens kernel $\{$diagonal, crossed$\}$). Aggregation is
exactly the collapsed $\Rrec = \Kbw$ interface with no forced joint
self-coordinate. The subject boundary is gate-determined by the surviving
promoted carrier, not by a topological surface or by summing micro-packages.
Interface contrast on the same states: the subject interface hides pairing
(exposure false); the collective-ledger interface publishes it (exposure
true), failing C4 by design, an institution-shaped control.
\emph{Nonclaim:} no general verdict about real collectives, and aggregation
of local contents is never sufficient: it is the collapsed regime.
\end{theorem}

\paragraph{T1b: the taxonomy, briefly.}
The classifier lattice of Section~\ref{sec:arch} belongs to this cluster
operationally: its four cells (controller, bound-without-drive,
drive-without-binding, full subject) are assembled from independently
witnessed T7 and T5 axis values, and it is a classification schema over
degraded packages, not four integrated carriers and not an enumeration of
all possible minds.

\begin{table}[p]
\centering\footnotesize
\begin{tabular}{p{0.05\linewidth}p{0.22\linewidth}p{0.3\linewidth}p{0.16\linewidth}p{0.15\linewidth}}
\toprule
Law & Mathematical type & Central exact result & Interface / assumption & Lean scope \\
\midrule
T1 & forcing + minimal refinement & $\Kviab(q)=\emptyset$ exhaustively; $M$ restores $10$ states; unique minimal $7$-class repair & deadline; live hidden phase & carrier + quotient tables; searches Python-certified \\
T1b & classifier schema & four cells from independent C7/C5 axes & landed axis witnesses & classifier + role table \\
T2 & non-exposure dichotomy & loop moves $M$ ($1/1$), fixes $Q$ ($0/1$); dissolution control & declared channel family & core facts; fractions pinned \\
T3 & nonfactorization corollary & strict pair $\Rightarrow$ no factor map; control factors & declared source/target maps & corollary on T9 support \\
T4 & quotient + certificate contrast & public match exhaustive; certificate battery diverges; C publicly discernible & declared public lenses; imported battery & records mirrored; enumeration Python-certified \\
T5 & priced residue + credit audit & scores $1,1,1/2$; products $8/27$ vs $1/27$ (self), $0$ vs $16/81$ (other), balanced (none); DPI null & bounded budgets; declared provenance & lift + credit rule owned; constants pinned \\
T6 & self-audit limits & own-level \texttt{undefined\_circular} $\times 3$; bridges provisional; confabulation completion & instrument-indexed claims & vendored abstract tower + toy statuses \\
T7 & strict recombination & $\Kbw{=}(2)$, $\Rrec{=}(1,1)$; marking collapses, erasure restores & declared branch family & criterion + witnesses \\
T8 & funded clock + holonomy & progress $2/3$, tick $2/9$, $\Delta A{=}2$; holonomy $1$ vs $0$; flow-null blocked & declared clock/protocol/budget & traces + noncommutation; rates pinned \\
T9 & report descent & non-injective report; no-report channel separates; overread blocked & declared report interfaces & witness + classifier \\
T10 & boundary schema & up $[0,0,1,2,3]$ / down $[3,3,2,1,0]$; thresholds $1/2$ vs $1/4$ & toy maintenance rule and grid & sweep constants pinned \\
T11 & exact LP duality & $\lambda=(24/11,14/11)$; slack $0$; curve monotone; proxy unstable & declared LP; valence is \gradeR & LP values pinned + relations \\
T12 & induced-kernel geometry & $43/270$ vs $1/810$; defect $13441/52488$; signature residues $0/1$ vs $1/1$ & declared lens/prototypes; corrected values only & orders/ratios via cross-multiplication \\
T13 & gauge rigidification & $G \cong \mathbb{Z}_2 \to \{\mathrm{id}\}$; transports $2 \to 1$ & declared descents; modeled rigidifier & counts pinned \\
T14 & lineage transfer & trauma $0/1$ vs $1/1$; fission/fusion; empty target & declared route package & lineage facts; death extension Python-only \\
T15 & combination gates & $\Rrec \succ \Kbw$ + refinement $1\to2\to2$ on one joint carrier & two interfaces, one carrier & stability/equivalence; refinement Python-executed \\
T16 & convention-relative composition & flags $(\top,\top,\top)$ $\to$ hard-problem report; control $\to$ certified & declared completion convention & composite facts + both completions \\
T18 & temporal cluster & strictness witness; defects $1{\to}0$, $4/7{\to}0$; entrainment $1/1$, $1/3$; transports $2/1$; T18c demoted & T8 base; T9 discipline & cluster constants pinned \\
\bottomrule
\end{tabular}
\caption{\textbf{Tbl-4: the consequence catalog, compactly.} One row per law
with its exact decisive result and mechanization scope; T17 has its own table
in Section~\ref{sec:t17}. Full six-field entries (setup, result, witness,
controls, grade, nonclaims) are in Appendix~\ref{app:compendium}.
Scope/price: every row is carrier- and interface-relative; ``Lean scope''
states what is proved there versus pinned as certified constants.}
\label{tbl:catalog}
\end{table}

 \section{Conditional realization and the open integrated target}
\label{sec:t17}


Everything so far is Tier~S. This section is where the ontological price is
paid: visibly, once, and only here. The realization umbrella (repository
name T17) is \emph{seven separately graded objects}, each with one
authoritative statement and one exact grade; a machine-readable claim
registry in the artifact pins the seven-way contract and is validated against
the Lean, Python, and prose surfaces. Nothing below is ever summarized as
``one realization theorem.''

\subsection{The two postulates}

\begin{postulate}[A1: Ontic Closure Completeness]
\label{post:a1}
To be real is to be a survivor of audited closure emergence, a lawful
promoted package that passes its gates, with no further realness
ingredient.
\end{postulate}

\begin{postulate}[A2: Substrate Indifference]
\label{post:a2}
Closure laws apply identically across carriers; the carrier label is not an
extra closure law.
\end{postulate}

Both are philosophical postulates \gradeO, stated in the design layer and
defended in a companion philosophical paper; neither is proved by the
structural catalog or by Lean, and the paper's Lean development being
polymorphic in the carrier type is \emph{not} a proof of A2: A2 currently
has no carrier-transport or invariance theorem at all, and its only formal
footprint below is a parametric scope. A reader may refuse either postulate;
the refusal costs the reader Section~\ref{sec:t17} and nothing else.

\subsection{The seven claims}

\begin{description}
\item[T17-S (simplified certificate schema).] \gradeS{} structural schema.
Lean defines a carrier-indexed \emph{simplified} certificate
$\Certsimp(X)$: mechanics-based clause predicates C1--C8 over one declared
transition/action family and lens family (the C1/C5/C6/C8 audit objects are
equation-coupled to the declared mechanics; C4/C7 are structurally
constrained on the same support but not transition-coupled, a schema-grade
distinction recorded explicitly), cross-clause coherence requiring one
nonempty communicating viable hub containing the clause witnesses, and a
simplified promotion ledger. The schemas are not ordered by any landed
theorem: $\Certsimp$ uses a weaker imported gate ledger while adding one-step
restrictions absent from the full design definition; no implication or
equivalence between $\Certsimp$ and $\Cpkg$ has been proved or is claimed,
and no integrated instance is exhibited. T17-S does not claim its schema
suffices for lawful closure or consciousness; that is exactly the
additional purchase in T17-O.
\item[T17-P (witness catalog and hygiene protocol).] \gradeS{}
protocol/catalog. The Python protocol imports and re-verifies eight per-law
witness fields (one row per clause, sourced from the landed laws of
Sections~\ref{sec:t1}--\ref{sec:clusters}) and recursively scans stored
condition-bearing values for smuggled vocabulary. Every current row carries
the verdict \texttt{witness\_verified}; no row label is a realization
verdict, and the protocol constructs no candidate, no integrated carrier, no
gate-panel object, and no realization result.
\item[T17-M (simplified certificate contrast).] \gradeS{} schema. Lean
defines $\mathrm{SimplifiedCertificateContrast}(a,b) :=
\mathrm{PublicMatch}(a,b) \wedge \Certsimp(a) \wedge \neg\Certsimp(b)$:
dynamic public-statistics matching over the complete public observation, a
simplified certificate for $a$, and its absence for $b$. This is an
explicitly simplified approximation of T4b's demarcation target, not the
full-certificate relation, which remains open until a full certificate
exists.
\item[T17-O (conditional realization bridge).] \gradeO{} axiom-relative
conditional schema; the exact logical form is displayed below.
\item[T17-I (integrated instance).] \gradeOpen. The required deliverable is
absent: an exhibited carrier and package satisfying a \emph{full} design-level
certificate, together with a \emph{proved} (not assumed) certificate-soundness
theorem. Both artifacts record this explicitly: Lean pins
$\texttt{singleCarrierInstanceExhibited}=\texttt{false}$ and Python pins the
same flag with realization grade \texttt{schema}.
\item[T17a (report-interface pointer).] \gradeS{} interface-relative pointer.
On the finite T3/T9 carrier, hidden no-report content does not factor through
the declared report channel. This says nothing about whether a designer could
specify or predict content from the complete state, initial conditions, or
dynamics.
\item[T17b (declared-interface opacity pointer).] \gradeS{}
interface-relative pointer. Relative to the declared current channels, T2's
hidden structure does not descend and T9 blocks overreading the report
boundary. Opacity may disappear under an expanded interface (T2's own
dissolution control demonstrates this); this is not opacity to every creator
or instrument, and it carries no moral-status conclusion.
\end{description}

\subsection{The exact form of the bridge}

\begin{theorem}[Conditional realization bridge; T17-O]
\label{thm:t17o}
\gradeO{} \emph{Grade:} axiom-relative conditional schema, mechanized.
\emph{Scope:} arbitrary carrier type $X$ and package $\mathrm{pkg}$;
$\mathrm{ClosureComplete}$ and $\mathrm{Realized}$ are caller-supplied
predicate parameters. Assume:
\begin{itemize}
\item[(H1)] \emph{certificate soundness}: every simplified certificate is
sound for lawful audited closure,
$\Certsimp(X, \mathrm{pkg}) \Rightarrow \mathrm{ClosureComplete}(X,
\mathrm{pkg})$, an unproved schema premise that is neither A1 nor A2; and
\item[(H2)] \emph{A1}: lawful audited closure is realization,
$\mathrm{ClosureComplete}(X, \mathrm{pkg}) \Rightarrow
\mathrm{Realized}(X, \mathrm{pkg})$.
\end{itemize}
Then every simplified-certified package is realized:
$\Certsimp(X, \mathrm{pkg}) \Rightarrow \mathrm{Realized}(X, \mathrm{pkg})$,
uniformly in the carrier type. \emph{Nonclaim:} neither premise is proved;
no A2 transport or invariance theorem exists (the universal carrier parameter
is A2's only footprint, not a proof step); no package is realized
unconditionally; and no candidate is certified.
\end{theorem}

The theorem's mathematical content is thin by design: a
composition of two assumed implications, checked by machine so that its
premises cannot silently grow or shrink. Its scientific content is the
\emph{pricing}: it fixes exactly what a reader must accept for the structural
catalog to acquire ontological weight, and it separates the two purchases so
they can be refused independently. Discrepancies between the generic
simplified clauses and the concrete source constructions (the boundaries
flagged in Section~\ref{sec:fullvssimp}) price H1, not the implication.

\subsection{The negative-attack corpus}

The simplified schema earned its current form adversarially: the artifact
records more than two dozen executable blocked attacks. Each is a concrete
package that a previous version of some clause wrongly accepted or would
wrongly accept, now rejected by a named Lean fact, frequently with a paired
predecessor fact recording the genuine old-pass/new-fail surface. Recurring
attack families: disconnected packages whose clause witnesses live outside
any communicating viable hub; hiddenness confined to underfunded states;
joint decoders that expose a determining state two channels hide
individually; predictive splits supported only by infeasible actions; merged
classes with divergent enabled actions; zero-cost (unpriced) existence;
manufactured credit symmetry from duplicated action lists; and
public-match spoofing via undeclared feasible actions or divergent channels.
Appendix~\ref{app:t17} tabulates the corpus. The methodological point is the
lane discipline again: a certificate schema that has survived named,
executable attacks is still only a schema. The attacks harden $\Certsimp$;
they do not inhabit it.

\begin{table}[t]
\centering\small
\begin{tabular}{p{0.09\linewidth}p{0.08\linewidth}p{0.3\linewidth}p{0.22\linewidth}p{0.17\linewidth}}
\toprule
Claim & Grade & Proved / executed object & Missing purchase & Status \\
\midrule
T17-S & \gradeS & carrier-indexed $\Certsimp$: clauses, coherence, simplified ledger (Lean) & no ordering vs.\ full $\Cpkg$; no instance & schema \\
T17-P & \gradeS & eight re-verified witness rows + hygiene scan (Python) & rows are not a certificate & catalog \\
T17-M & \gradeS & PublicMatch $\wedge$ $\Certsimp(a)$ $\wedge$ $\neg\Certsimp(b)$ (Lean) & full-certificate contrast & schema \\
T17-O & \gradeO & conditional bridge, uniform in carrier (Lean) & H1 soundness; A1; A2 unformalized & conditional only \\
T17-I & \gradeOpen & explicit \texttt{false} marker in both artifacts & the integrated instance itself & open \\
T17a & \gradeS & report-interface nonfactorization pointer & n/a & pointer \\
T17b & \gradeS & declared-interface opacity pointer & n/a & pointer \\
\bottomrule
\end{tabular}
\caption{\textbf{Tbl-5: the seven-way realization split.} Scope/price: each
row has one grade and one authoritative statement; the machine-readable
registry in the artifact enforces the split across Lean, Python, and prose.}
\label{tbl:t17}
\end{table}

\begin{figure}[t]
\centering
\fbox{\parbox{0.92\linewidth}{\small\emph{[Figure F6: realization ladder;
rendered per the artifact figure plan.]} T17-S, T17-P, and T17-M as separate
structural rungs; the T17-O implication drawn with its two premises
(certificate soundness; A1) as explicit unspent inputs and A2 shown only as a
carrier-parametric footprint, never as a proved transport arrow; the ladder
terminates at T17-I, labeled \textbf{OPEN} inside the figure.}}
\caption{\textbf{F6: the realization ladder.} Takeaway: structure, catalog,
contrast, conditional bridge, and open instance are five different things.
Scope restriction: an ontological premise never looks like a theorem arrow;
the eight witness rows are never merged visually into a certificate.}
\label{fig:ladder}
\end{figure}

\subsection{What refusal costs}

The architecture is built so that its ontological reading is refusable
without wreckage. Decline H1 and the structural catalog stands untouched:
one simply does not read $\Certsimp$ as sufficient for lawful closure.
Decline A1 and lawful closure remains a mathematical property: one simply
does not read it as realization. Decline A2 and every result remains
carrier-indexed. What cannot be done, on pain of the factorization error of
Section~\ref{sec:selfaudit}, is to accept the strict-promotion certificates
and then paraphrase them away; and what cannot be claimed, on pain of
misreading Table~\ref{tbl:t17}, is that anything in this paper certifies an
actual subject. It does not, and it says so in machine-checked form.

 \section{Empirical status and the prospective interface}
\label{sec:empirical}


This section exists to make one negative fact and its consequences precise.

\emph{No empirical data were collected, fitted, or used to validate the
architecture.} The catalog's evidence is exact finite computation; exact
finite studies are not experiments on consciousness, and nothing in
Sections~\ref{sec:t1}--\ref{sec:t17} is empirical evidence for anything about
any actual system.

\emph{The inherited literature battery is not a contribution.} An earlier
preparation phase assembled a post hoc battery of literature-derived
phenomenon candidates (perceptual, temporal, and dissociation phenomena) with
proposed clause mappings. Screened against its own admission protocol
(verified bibliographic metadata, a checked load-bearing passage, same-sample
altered \emph{and} preserved observables, an independently frozen mapping,
and a row-level defeat condition), the battery failed its row and
source-verification requirements. It is therefore excluded from this paper's
results and retained in the artifact only as non-evidential protocol history
(status \texttt{history\_only\_non\_evidential}, per-row;
Appendix~\ref{app:empirical}). The candidate set was assembled after the
architecture existed and is not unbiased, preregistered, or predictive of
anything; no phenomenon row appears in this paper as an adequacy result.

\emph{What a real test would require.} Table~\ref{tbl:obligations} states the
prospective obligations. Four ingredients are missing and are all mandatory.
An \emph{observable bridge}: an audited descent from each formal target to
a measured report/behavior record, in T9's sense, rather than a verbal
mapping. A \emph{prospective mapping freeze}: clause-to-observable mappings
fixed before data contact. An \emph{independent candidate universe}: a
selection rule for test phenomena that does not condition on the
architecture's successes. And a \emph{same-sample selective-dissociation
design}: within one sample, a manipulation predicted to alter the mapped
observable while a declared control observable is preserved, with the
dissociation's failure predeclared as the row's defeat condition. These
requirements are the practical face of the testability critiques that
motivate them \citep{KleinerHoel2021, DoerigEtAl2020HardCriteria}.

\emph{Two honest limits on falsifiability.} First, the present
parameterization does not support a sharp package-level empirical falsifier:
the architecture's exact content lives on declared toy carriers, and what a
failed prospective row would defeat is its row (mapping and clause reading),
not the mathematics. Structural countermodels (the executable attack
corpus of Section~\ref{sec:t17}) are not empirical falsification in the sense
formalized by \citet{KleinerHoel2021}. Second, no external C8 row is
admissible under the current scope: mapping the viability clause to any
human existence boundary would require clinical evidence types
(loss/recovery of consciousness, anesthesia depth) that the current
human-only, non-neural, report-level scope excludes; T10's hysteresis is a
property of a toy accumulator until then.

\begin{table}[t]
\centering\footnotesize
\begin{tabular}{p{0.14\linewidth}p{0.2\linewidth}p{0.24\linewidth}p{0.2\linewidth}p{0.08\linewidth}}
\toprule
Formal target & Missing observable bridge & Required prospective design & Predeclared defeat condition & Status \\
\midrule
C4/T2 hiddenness & descent from declared-channel non-exposure to a report/decoding record & fixed channel family; decoding attempt predicted to fail while a control decode succeeds & successful decode inside the declared family & not designed \\
C7/T7 binding & descent from $\Rrec \succ \Kbw$ to a same-sample conjunction measure & marking manipulation predicted to collapse a joint-content measure, preserving branchwise measures & no marking-induced collapse & not designed \\
C6/T8 temporal funding & descent from funded tick structure to a load-manipulated timing record & resource manipulation predicted to alter tick-like regularity, preserving order judgments & no funding sensitivity & not designed \\
D1/T11 price profile & descent from binding-constraint duals to a graded motivational record & constraint tightening predicted to shift a graded measure monotonically; slack regime flat & non-monotone or slack-sensitive response & not designed \\
C8/T10 boundary & \emph{none admissible}: external C8 mapping excluded under current scope & clinical loss/recovery designs out of scope & n/a & excluded \\
\bottomrule
\end{tabular}
\caption{\textbf{Tbl-6: prospective empirical obligations.} These are
obligations for future work, not results, and the rows are illustrative
targets: every field must be frozen prospectively before any data contact.
The inherited candidate battery supplies \emph{no} rows here; it is
non-evidential history (Appendix~\ref{app:empirical}). Scope/price: a failed
row defeats its mapping row, not the Tier~S mathematics; there is currently
no sharp package-level falsifier.}
\label{tbl:obligations}
\end{table}

 \section{Related work and demarcation}
\label{sec:related}


The method of this section is uniform: state each comparator in its strongest
relevant form, state the exact overlap, state the exact difference, and claim
nothing about a comparator that its own definitions do not support. No
comparison here asserts that a neighboring theory is refuted by an omitted
clause, and no ``every previous theory found one clause'' claim is made.

\subsection{Mathematical consciousness science and formal theory testing}

The paper's home field is mathematical consciousness science.
\citet{Kleiner2020MathematicalModels} gives a general framework for what a
mathematical model of consciousness must declare (experience spaces, a
physical domain, and the mappings between them) and criteria for evaluating
such models. \citet{KleinerTull2021} axiomatize generalized IIT as a
mathematical structure: system classes, experience spaces, decompositions,
integration levels, cores, and a system-state-to-experience map; describing
that work as chiefly category-theoretic would undersell it, since it is a
formalization of a theory class. \citet{SignorelliWangCoecke2021} develop a
process-theoretic, compositional treatment of conscious experience and
already derive formal results about phenomenal unity and a form of
privacy/unreadability within their calculus. The present paper works inside
the field these papers establish and adds a different kind of object: an
audit-contract architecture with exhaustively computed finite carriers,
per-sentence claim grading, and an explicitly priced ontological bridge. The
overlap with \citet{SignorelliWangCoecke2021} deserves the sharpest
statement: they formalize unity and unreadability compositionally; T2 and T7
study public-interface factorization and strict recombination refinement on
finite carriers with executable dissolution/collapse controls. These are
different formal mechanisms for adjacent explananda, and this paper claims no
priority for formal unity or privacy results.

On testing: \citet{KleinerHoel2021} formalize when inference procedures make
consciousness theories falsification-resistant or trivially falsified;
\citet{DoerigEtAl2020HardCriteria} state hard criteria an empirical theory
should meet; \citet{CogitateConsortium2025} ran a preregistered adversarial
collaboration that challenged central predictions of both IIT and global
workspace theory without producing a definitive verdict. This paper's
relation to that literature is candid: it does not
satisfy the hard criteria, is not tested by COGITATE-style designs, and says
so (Section~\ref{sec:empirical}); its contribution to the testing discussion
is the obligations table and the observation that structural countermodels
are a different object from empirical falsification.

\subsection{Information Closure Theory}

The closest named closure comparator is Information Closure Theory
\citep{ChangEtAl2020InformationClosure}, which defines conscious processes as
coarse-grained processes with positive nontrivial informational closure
(NTIC): the macro process predicts its own future as well as the environment
does, without that prediction routing through trivial self-copying. The
shared conviction is that closure, a description layer that accounts for
its own continuation, is the right kind of property to make precise. The
formal objects differ throughout: ICT's conditions are
information-theoretic (mutual-information quantities over process
trajectories); C1--C8 are structural (closure operators, factorization
defects, refinement terminality, viability kernels, audit gates). T1 in
particular computes viability-relevant predictive and refinement facts, not
NTIC, and no result in this paper computes ICT's quantities. Accordingly, no
reduction and no non-equivalence theorem between the formalisms is claimed in
either direction; the comparison is between two independently defined notions
of closure, one informational and one audited-structural.

\subsection{Predictive states, viability, and regulation}

Section~\ref{sec:t1position} already gave the load-bearing demarcations, so
one summary suffices. The predictive quotient is the finite controlled
analog of causal states \citep{CrutchfieldYoung1989, ShaliziCrutchfield2001};
what is new in T1 is the question (does a bounded public self-lens
\emph{survive}), the exhaustive answer (no, on the declared carrier), and the
minimal-repair object (a refinement of a declared lens, which causal-state
minimality does not address). The viability kernel and its iteration are
standard \citep{AnayaEtAl2009Viability, AubinBayenSaintPierre2011}; the
paper's use makes kernels properties of interfaces. The good-regulator
theorem \citep{ConantAshby1970} is retained only as motivation and with its
own narrow content stated.

\subsection{Integrated Information Theory 4.0}

IIT 4.0 \citep{AlbantakisEtAl2023IIT4} is the field's most developed
axiomatic architecture: axioms about experience converted into postulates
about intrinsic cause-effect structure, with integrated-information
quantities, distinctions and relations, and an exclusion postulate selecting
maximal complexes. The structural ambitions overlap: both programs posit
that being a subject is a formal property of a structure, and both take
quality structure seriously (IIT through cause-effect structures, this paper
through descended similarity geometry and gauge). The differences are
categorical rather than quantitative. First, no quantity: this paper computes
no integrated information, no cause-effect structure, and no
$\Phi$-like scalar, and its finite structural vocabulary must not be read as
IIT terminology. Second, identity versus certification: IIT asserts an
identity between experience and maximal cause-effect structure; this paper's
strongest ontological move is a conditional bridge with two refusable
premises, and its certificate schema is explicitly not a consciousness
detector. Third, exclusion versus gating: where IIT's exclusion postulate
selects one complex, T15's combination law gates joint subjects by two
structural certificates on one carrier and leaves ``which packages exist'' to
viability rather than maximality. This paper neither tests nor refutes IIT.

\subsection{Quality spaces, structuralism, and gauge}

Structuralist programs hold that the scientifically tractable content of
experience is relational structure: quality spaces with their similarity
geometry. Work in this tradition includes principled accounts of spatial
experience structure \citep{HaunTononi2019}, category-theoretic and
relational approaches \citep{TsuchiyaSaigo2021}, and topologically structured
phenomenal spaces \citep{Prentner2019Topological}. T12 and T13 are this
paper's contribution to that conversation, and their boundary is stated from
this side: T12's landed object is an exact descended finite similarity
geometry with coherence audits on a declared carrier, not an empirical
phenomenal space, not a psychophysical fit, and not a general geometry-birth
theorem; T13 turns inversion talk into a computable gauge group and shows
exactly one descended asymmetry collapsing it. No clause-level comparison
with the cited programs is made beyond this, and none of them is claimed to
be superseded. On the philosophy side, the paper's debts are explicit:
subjective character as explanandum \citep{Nagel1974}, the explanatory-gap
label \citep{Levine1983}, the phenomenal/access distinction
\citep{Block1995}, and the hard-problem and meta-problem framings
\citep{Chalmers1995FacingUp, Chalmers2018MetaProblem}. These are explananda
and distinctions the architecture addresses; they are premises of nothing in
the catalog.

Substrate questions are live in the adjacent literature, for example a
recent no-go argument concerning consciousness on a chip
\citep{KleinerLudwig2024Neurons}. This paper's A2 is a postulate with no
formal transport theorem (Section~\ref{sec:t17}); it takes no position on
such no-go arguments and cannot, until a carrier-transport theorem exists on
its own side.

\subsection{Six Birds provenance and owned seams}

The imported machinery (promotion bridges, factorization defects,
no-overreading, same-level self-audit limits, rotating audits, hiddenness
normal forms, route transport, recombination quotients, protocol traps,
closure deficits, clocks, geometry pipelines, and currency duality) comes from
the companion corpus cataloged in Appendix~\ref{app:imports}, cited there
theorem by theorem with hypotheses. Sibling citations show provenance and
ownership, never independent corroboration. Three adjacent Six Birds papers
own seams that this paper does not cross
\citep{Tsiokos2026FoundationsVEndogenousClosure,
Tsiokos2026FoundationsVIDynamicalLaws,
Tsiokos2026InstitutionsStrictExtensions}. The cognition paper (Foundations~V)
owns operational selfhood, individuation, priced attention allocation,
control-price compression, bounded-reflexivity capacity economics, and repair
transport; this paper's T6 owns only its finite audit witness, its D3 is
recognition-grade, and no claim here fills Foundations~V's deferred
global-access slot. The dynamical-laws paper (Foundations~VI) owns
theorist-side properties of unbounded runs (solvency, escrow, carry-horizon
confinement); this paper's clocks and rewrites are finite package-internal
certificates, and a funded tick is neither eventual liveness nor a
telescoping cost bound. The institutions paper owns its repeated-game
promotion and formation theorems; T1 shares an obligation skeleton with them
on a different carrier (carrier non-definability, saturation, strictness),
and no result is transported: in particular, T1 computes no
closure-deficit quantity and its canonical nine-class quotient is never
identified with its seven-class minimal refinement. No retained sibling
proposition is a logical premise of any theorem in this paper.

 \section{Limitations, open program, and conclusion}
\label{sec:limits}


\subsection{Limitations}
\label{sec:limitations}

Twelve limitations, before any conclusion.

\begin{enumerate}
\item \textbf{Exact finite models only.} Every positive result lives on a
declared toy carrier. Nothing licenses transfer to large, noisy, or
continuous systems beyond the displayed conditional schemas.
\item \textbf{No empirical content.} No dataset, no parameter fit, no
validation, no phenomenon rows (Section~\ref{sec:empirical}).
\item \textbf{Provenance of predictions.} The artifact contains documents
labelled as predictions, but their initial Git history does not establish
temporal precedence over implementations, and later corrections are
versioned. Nothing in this paper calls them preregistered.
\item \textbf{Selected mechanization.} Lean coverage is selected, not
catalog-wide: kernels prove core facts and pin many enumerated values as
certified constants matching the Python computations (per-law boundaries in
Table~\ref{tbl:catalog} and Appendix~\ref{app:repro}); the current
manifest-level validator accepts a zero-row schema and is a known release
blocker, so no manifest-audited coverage claim is made.
\item \textbf{Inherited trust base.} The Lean development trusts vendored
foundations modules and declared axioms from the companion corpus, cataloged
in Appendix~\ref{app:repro}; their opacity guards are inherited, not
re-proved.
\item \textbf{T1's family boundary.} The forcing theorem is exact on one
carrier; the family schema's hypotheses (live non-definable phase, deadline
relevance) are strong and not claimed common (Section~\ref{sec:t1family}).
\item \textbf{T12's open target.} Geometry birth (fitted distortion,
dimension, curvature transport) is not landed; the landed object is an
audited induced geometry.
\item \textbf{T16's conventions.} The meta-problem composition is relative
to a declared completion convention, and its carrier collapses several
premises extensionally.
\item \textbf{T17's open core.} The simplified/full schema relation is
unproved in both directions, certificate soundness is assumed where used,
A2 has no formal theorem, and T17-I, the integrated instance, is open
(Section~\ref{sec:t17}).
\item \textbf{T18c demoted.} The temporal-precondition claim is a dependency
observation, not a theorem.
\item \textbf{Recognition bridges are unproved.} Valence, phenomenological
clause mappings, attention readings, and all state-taxonomy readings are
named hypotheses (\gradeR), each retractable without structural loss.
\item \textbf{Review status.} The artifact's ``external review request''
documents record an internal author-driven review loop; they are not
independent peer review, and this preprint has not been peer reviewed.
\end{enumerate}

\subsection{Open program}

The open obligations, in order of weight: exhibit T17-I or prove that some
declared obstruction prevents it; prove or refute an ordering between
$\Certsimp$ and the full $\Cpkg$; formalize A2 as a carrier-transport theorem
or retire it to purely philosophical status; extend the simplified C3 from
one-step stability to a terminal bisimulation; land the T12 geometry-birth
checklist; strengthen T18c or retire it; and design one prospective empirical
row to the standard of Table~\ref{tbl:obligations}.

\subsection{Conclusion}

This paper's contributions, restated once in their fixed order: a claim
architecture that keeps structural results, recognition bridges, ontological
implications, and open targets in separate, artifact-enforced lanes; an
eight-clause audited closure contract kept explicitly distinct from its
mechanized simplification and from any instance; one exact forcing result,
on whose fully declared deadline carrier every public-lens policy fails,
the canonical predictive quotient survives, and the unique minimal repair is a
seven-class predictive split of one public fiber; a catalog of exact,
interface-relative structural recastings of the classical problems; a
seven-way graded realization program whose ontological bridge is mechanized,
priced, and refusable, with its integrated target visibly open; and an
audited artifact whose remaining gaps are listed rather than implied away.

The strongest exact result is Theorem~\ref{thm:t1}. The strongest open
obligation is T17-I. Between them lies the paper's actual thesis, which is
double and must remain double: structurally, a declared closure package
organizes a family of exact finite results about self-prediction, privacy,
report, binding, temporality, geometry, identity, combination, and audit
limits; ontologically, realization follows only as a conditional bridge after
purchasing certificate soundness and Ontic Closure Completeness, and the
integrated instance that would give that bridge something unconditional to
say does not yet exist. The two halves of the thesis are printed in different
lanes, and the paper ends where its evidence ends.
 
\appendix
\section{Definitions and imported machinery}
\label{app:imports}


\subsection{Notation summary}

\begin{center}\small
\begin{tabular}{ll}
\toprule
Symbol & Object \\
\midrule
$X$ (T1: $Z = P \times W$) & finite carrier \\
$A$, $F_a$ / $P_a$ & actions; deterministic maps / stochastic kernels \\
$q : X \to Y$; $\Pi_q$ & public lens; its observational partition \\
$U$ & hidden/promoted (non-descending) object \\
$Q$, $M$, $\pi : M \to Q$ & current public and predictive quotients; comparison map \\
$L$ & T1's continuation-law abstraction (never renamed $M$) \\
$\Pi_*$ & T1's unique minimal closure-restoring refinement (7 classes) \\
$\Dfact(\pi_0,\pi_1)$ & factorization defect (set of split pairs) \\
$\Kviab(\ell)$ & viability kernel under $\ell$-measurable policies \\
$\Kbw$, $\Rrec$, $\eta:\Rrec\to\Kbw$ & branchwise and recombination quotients; forgetful map \\
$E_\tau$; $\delta$ & staged closure operator; idempotence defect \\
$\lambda_i = \partial\Phi^*/\partial b_i$ & dual price of budget $b_i$ \\
$\Cpkg$; $\Certsimp(X)$ & full C-package target; simplified certificate schema \\
$\mathrm{Realized}$, $\mathrm{ClosureComplete}$ & caller-supplied predicates in T17-O \\
\bottomrule
\end{tabular}
\end{center}

Fixed distinctions the notation enforces: $Q$, $L$, and $M$ are three objects
in T1; the nine-class canonical $M$, the seven-class full-carrier $\Pi_*$, and
the seven-class kernel-restricted terminal quotient are pairwise distinct;
$\Kviab$ and $\Kbw$ are unrelated despite the inherited letter; public match,
hidden-state difference, certificate contrast, and realization are four
things; structural dual price and felt valence are two things; and agency,
cognition, attention, global access, report, and consciousness are
non-equivalent throughout.

\subsection{A minimal corpus primer}

The companion corpus (``Six Birds'') is a discipline for \emph{audited
emergence}: when does a coarse description become a theory with stable
objects, how does strict extension of what is definable remain possible, and
how do directionality and self-audit claims stay honest under observation.
Its first foundations paper isolates the operator toolkit (theory packages
$T = (Z, f, \Sigma_f, E, A)$: carrier, lens, definability algebra, completion
endomap, audit functional; closure operators and ladders; drive/affinity
audits). The second fixes the admissibility meta-theory over finite audited
typed closure descriptions, including the channel-versus-activation
distinction (six role \emph{channels} exist; nothing requires six
\emph{active} roles). The third builds the finite audited interaction
calculus this paper leans on hardest: typed cells, instruments,
promotion bridges with gate panels, defect families, and the self-audit
theorems. The fourth is a catalog of layer-agnostic normal-form laws
(hiddenness, irreducibility, and others) with an explicit
substrate-pending/obligation grading. Around the foundations sit
phenomenon papers, each owning one engine imported below. All are cited here
as corpus-internal sources with pinned versions in the artifact's import
inventory; canonical public locators are pending owner confirmation and will
be added before submission.

\subsection{Imported theorem map}

Each import is used under its own hypotheses; analogy is never theorem
transport. The table names the source engine, the imported object, and the
consuming law.

\begin{center}\footnotesize
\begin{tabular}{p{0.23\linewidth}p{0.42\linewidth}p{0.24\linewidth}}
\toprule
Source (corpus paper) & Imported object & Used by \\
\midrule
Foundations I & counting lemma ($2^K$ of $2^N$ predicates definable; strict-extension corollary) & T1 genericity framing \\
Foundations II & channel-vs-activation discipline; recognition-hypothesis threat & T1b recognition audit \\
Foundations III & nonfactorization dichotomy (Thms 12/13); promotion-gate soundness (Thm 19); strict promotion $\Rightarrow$ nonfactorization (Thm 20); no-overreading (Thm 23); same-level self-audit failure (Thm 24); finite rotating audit (Thm 25); gate/status vocabulary & T3, T5, T6, T9, T16, T17-S ledger \\
Foundations IV & hiddenness normal form (F13a) and formed-closure law (F13b), substrate-pending/obligation grades; irreducibility as no-shorter-closed-description (F42) & T2 framing, T3 \\
Hiddenness paper (HIDD) & core structural law (dissolution/preservation/degenerate/strength-independence); exposure predicate & T2, T17-S C4 \\
Route-transport paper (HOLMEM) & $Q$/$M$ quotients, comparison map, predictive transport, loop asymmetry, reachable-image factorization & T1, T13, T14, T18e \\
Recombination paper (MER) & branchwise/recombination quotients $\Kbw$, $\Rrec$, forgetful map $\eta$, strict-refinement witness criterion & T7, T15, T17-S C7 \\
Protocol-trap paper (PTRAP) & autonomous two-component lift; zero-EPR trap; arrow decomposition & T5 authorship, T8 \\
Closure-deficit paper (CAST) & closure deficit $\mathrm{CD}_\tau$; budgeted randomness curve & T5 residue \\
Time paper (NOTCH) & order/tick/arrow triple; purchased clock; anti-stall metrics; no-global-time holonomy; enablement- vs causation-time & T8, T18a/b \\
Geometry paper (PLOT) & induced macro kernel $\hat P = U P^\tau C$; path-cost metric; coherence audits & T12 \\
Currency paper (SPEND) & currency roles; shadow-price duality; proxy discipline & T11 \\
Agents paper (THROW) & viability kernel as greatest fixed point; agenthood/agency firewall; empowerment nulls & T1, T10, T11, T14 \\
Usefulness paper (WMEW) & audited-utility certificate soundness: an accepted certificate licenses the audited lower answer, never literalization of the hidden object & T9, T3 \\
Quantum paper (QSB) & empirical-indiscernibility principle; Leibniz-quotient factorization & T4a, T13 \\
No-go paper (NOGO) & deterministic data-processing inequality for arrow audits (no fake arrow by coarse observation) & T5, T8 nulls \\
Reflexive paper (REFLEX) & reflexive audit of packaging artifacts; open capacity-contraction conjecture (F3) & T6 boundary \\
Anti-reductionism primer & no-shortcut law; within/between-layer discipline; factorization error & T3, T5(vi), T16 \\
\bottomrule
\end{tabular}
\end{center}

Nonclaims inherited with the imports, stated once: the interaction calculus
proves \emph{no} general bridge-stacking composition and \emph{no}
local-to-global promotion theorem, so no law here composes bridges or glues
local certificates globally; hiddenness normal forms are
substrate-pending/obligation grade at the catalog layer; and no imported
result loses its own carrier restrictions by being used here.

\subsection{Trust base}

The Lean development vendors selected foundations modules (the interaction
calculus classifier, rotating-audit structure, and gate vocabulary) and
trusts their declared axioms, which are opacity-guarded in the source corpus
(axiom-applied wrappers guarded by opaque predicates with no public
constructor). The consciousness-side law kernels add no axioms of their own;
their per-declaration axiom dependencies are audited by
\texttt{\#print axioms} tooling recorded in the artifact. What Lean does and
does not establish per law is stated in Table~\ref{tbl:catalog} and
Appendix~\ref{app:repro}.

 \section{T1: exact carrier and proof details}
\label{app:t1}


\subsection{Carrier and quotients}

$Z = P \times W$, $|P| = 6$, $|W| = 2$; public lens $q(p,w) = p$; live phase
slip with exact branch weights $3/4$ and $1/4$; deadline making the hidden
continuation law viability-relevant. Public-history ambiguity witness: the
trace $(1,2,1)$ ends in both phases. Finite certificate: $N = 12$, $K = 6$,
$64$ definable predicates of $4096$ total ($4032$ non-definable); split
witness $p_1w_0 \to p_2w_1 \to p_1w_0$ versus $p_1w_0 \to p_2w_1 \to p_1w_1$;
strict witness $(p_1w_0, p_1w_1)$. The continuation-law quotient $L$ has five
classes implementing (need-cw $\to$ cw, need-ccw $\to$ ccw, return $\to$
stay); $L$ is incomparable with $q$ because its transit, lethal, and return
classes cross public fibers. Canonical $M = q \times L$: nine classes,
strictly refining $q$, splitting exactly the public fibers at positions
$1$, $2$, $5$.

\subsection{Policy and refinement enumeration}

$\Kviab(q) = \emptyset$ over every $q$-measurable feedback policy
(exhaustive); the stored $L$-policy kernel has ten states; the canonical
$M$ search returns the same ten states, with $\Kviab(L) = \Kviab(M)$
certified by the refinement audit. Single-fiber split kernel sizes across
$p = 0,\dots,5$: $(0, 10, 0, 0, 0, 0)$. The maximal refinement not splitting
$p{=}1$ has empty kernel; monotonicity of kernels under lens refinement (a
finer lens implements every coarser-lens policy) then closes the
$64$-refinement lattice: the unique minimal restoring refinement is the
seven-class $\Pi_*$ splitting exactly $p{=}1$. A separate, earlier 15-partition
audit over four reachable states remains an exploratory CORE-local audit
(exactly two CORE partitions survive, both factoring through $L$); it is not
part of the minimality proof. Controls: no-deadline gives
$\Kviab(q) \neq \emptyset$ with policy \emph{stay}; the washout/no-slip
control makes phase publicly reconstructible and is refused as forcing.

\subsection{Closure objecthood (C1)}

On the ten-state kernel under the canonical $M$ policy: closed-loop matrix
$P$; class indicator $C$; uniform prototype $U$ on realized $M$-fibers;
$E_\tau = P^\tau C U$. Multi-state transit fibers are transient under the
closed loop, so within-fiber prototype choice cannot change unperturbed
$E_\tau$; the declared $C$, $U$ content is pinned extensionally, and every
stored operator entry is checked against a fresh reconstruction from
$P, C, U$. Horizon table (rank/fixed-space dimension; defect): $\tau = 1$:
$10/0$, $\delta = 1$, unique CORE-uniform fixed point; $\tau = 2$: $9/1$,
$\delta = 1/8$, two-extremal fixed segment; $\tau = 3$: $10/0$, $\delta = 1$;
$\tau = 4$: $9/1$, $\delta = 3/32$. The even-horizon segment midpoint is the
odd-horizon fixed point. Perturbation at $\tau = 2$,
$P_\varepsilon = (1-\varepsilon)P + \varepsilon\,\mathrm{Unif}(K)$,
$\varepsilon = 1/16$: rank rises to $10$; the unique fixed point has mass
$38/155$ on each $p_1$ state, $1171/4960$ on each return state, and $1/160$
on each of the six transient states; CORE retention $77/80 \ge 3/4$; total
variation from CORE-uniform $3/80$; defect $78795/524288$. Second row:
$\varepsilon = 1/100$ gives total variation $3/500$; both rows satisfy
$\mathrm{TV} = (3/5)\varepsilon$.

\subsection{Terminal self-lens (C3)}

Profile refinement on the ten-state kernel from the six-class public
partition, splitting by exact per-action transition profiles with
out-of-kernel exits labeled distinctly: size trajectory $6 \to 7 \to 7$;
instability witness $(p_1w_0, p_1w_1)$; all $100$ ordered-pair equivalence
checks identify the terminal partition with canonical $M$ restricted to the
kernel. Extensionally $M(z) = (q(z), L(z))$.

\subsection{Lean mapping}

The kernel declaration \texttt{forcingOfI\_T1\_kernel} formalizes the
twelve-state carrier and the public, continuation-law, and canonical quotient
tables. Python-certified (with tests), not Lean-derived: the exhaustive gfp
policy searches, the $64$-refinement enumeration (whose monotonicity step is
a prose lemma over the eight computed searches), the $\tau = 2..4$ closure
rows, and the perturbation audit. Lean's C1 facts cover the $\tau = 1$
operator by evaluation with the defect witness; Lean's C3 facts prove
terminal-partition stability and public instability with $6/7/7$ as declared
constants.

 \section{Exact finite result compendium}
\label{app:compendium}


One six-field entry per law: \emph{setup / central result / decisive witness
/ controls / grade / nonclaims}. All values are the current corrected
snapshot values; superseded pre-correction values (notably in T12, T15, T16)
never appear. T1 is in Appendix~\ref{app:t1}; T17 in Appendix~\ref{app:t17}.

\paragraph{T1b (taxonomy and recognition).}
\emph{Setup:} T1's forcing cascade; contingent axes C7 (binding) and C5
(drive). \emph{Result:} four-cell classifier lattice (controller /
sentient-without-self / self-without-binding / full subject); scoped
recognition audit mapping C1--C8 to declared role composites with no seventh
role. \emph{Witness:} each axis value independently witnessed (T7 coupled
witness and aggregation collapse; T5 genuine-drive credit and internalized
zero-affinity). \emph{Controls:} drive-absent cell uses the scheduled/fake
control; binding-absent uses aggregation. \emph{Grade:} schema-grade cascade;
proposition-grade recognition audit. \emph{Nonclaims:} not four integrated
carriers; not the only possible structures; no universal six-role
completeness.

\paragraph{T2 (interface-relative hiddenness).}
\emph{Setup:} two hidden states $m_0, m_1$; declared channel family with
one-class current quotient $Q$; two-class predictive $M$. \emph{Result:}
connector loop (swap) fixes $Q$ and moves $M$; exposure false.
\emph{Witness:} moved-class fractions $0/1$ on $Q$, $1/1$ on $M$; declared
current observable computed constant on the hidden support. \emph{Controls:}
adding an $M$-reading channel dissolves the witness; the degenerate $M = Q$
case licenses nothing. \emph{Grade:} theorem-grade direct instantiation; C4
generic/declared, not universal. \emph{Nonclaims:} no unconditional privacy;
no metaphysical inaccessibility; hiddenness does not imply phenomenality.
Lean note: core facts proved on the wheel; moved-fraction numerators pinned
as certified constants.

\paragraph{T3 (no-shortcut corollary).}
\emph{Setup:} T9's support; report boundary as lower readout, no-report
boundary as target. \emph{Result:} strictness (nonempty $\Dfact$) $\iff$ no
factor map; no closed lower compression recovers strict content.
\emph{Witness:} factor attempt returns none on the strict case.
\emph{Controls:} report-to-report target factors concretely
($r_\mathrm{shared} \mapsto r_\mathrm{shared}$, $r_\mathrm{other} \mapsto
r_\mathrm{other}$); no gap certified without strictness. \emph{Grade:}
corollary-grade. \emph{Nonclaims:} no Tier~O assertion; mechanism not denied;
the philosophical gap not solved.

\paragraph{T4 (zombie study).}
\emph{Setup:} T2's public carrier face; imported certificate battery
(T2/T7/T5/T6). \emph{Result:} (a) layer quantities factor through the
Leibniz quotient; (b) A/B are one K-matched public face with genuine versus
degenerate certificate readings. \emph{Witness:} exhaustive public
comparison: identical multisets $\{\texttt{current}: (q,q)\}$; certificate
divergence on all four imported axes. \emph{Controls:} carrier C declares an
$m$-reader observable and is classified publicly discernible before
certificate comparison. \emph{Grade:} (a) theorem-grade; (b) schema-grade;
(c) deferred \gradeO{} pointer. \emph{Nonclaims:} layer-relative sameness
only; no integrated two-carrier pair; no zombie detector.

\paragraph{T5 (priced deliberation).}
\emph{Setup:} one deterministic support, records (public, last action,
hidden), choice $=$ hidden XOR public; budgets $B = 0,1,2$; declared
generator provenance. \emph{Result:} Theorem~\ref{thm:t5}'s six clauses.
\emph{Witness:} scores $1, 1, 1/2$; self products $8/27$ vs $1/27$; other
products $0$ vs $16/81$; internalized lift decomposes $P = C + D$ on all 81
entries, both channels symmetric, credited products $0$ vs $0$ (canonical
cycle $1/324$ each way); attribution triple (self, other, none).
\emph{Controls:} DPI null (coarse fluxes $1/27$ each way); out-of-scope and
bridge audit variants. \emph{Grade:} per clause (framing / theorem-grade toy
/ direct instantiation / direct instantiation / foreclosure / O-flavored
foreclosure). \emph{Nonclaims:} residue is bounded-observer ignorance;
classifier statuses are not hidden phenomenal reports; credit is
model-relative on declared channels. Lean note: lift, credit rule, and
attribution owned in Lean; support counts pinned.

\paragraph{T6 (introspection limits).}
\emph{Setup:} instrument-indexed self-audit claims; three-level tower;
completion toy with visible record and hidden true state. \emph{Result:}
same-level completeness blocked; towers non-final; unique closure-preserving
completions can confabulate. \emph{Witness:} statuses
(\texttt{undefined\_circular} $\times 3$); bridges provisional; missing-update
completion \emph{left} vs hidden \emph{right}. \emph{Controls:}
audit-complete case completes correctly; out-of-scope variant separates
status kinds. \emph{Grade:} direct instantiation plus finite extension.
\emph{Nonclaims:} introspection not impossible/worthless; F3 stability
explicitly open; no clinical claim. Lean note: vendored abstract
rotating-audit compatibility theorem has no Python counterpart by design;
Python owns the concrete tower statuses.

\paragraph{T7 (binding).}
\emph{Setup:} branch assemblages over modality/content branches; branchwise
$\Kbw$ vs recombination $\Rrec$. \emph{Result:} strict refinement certifies
binding; marking collapses it; erasure restores it. \emph{Witness:} coupled
regime $\Kbw = (2)$, $\Rrec = (1,1)$, pair (bound-diagonal, bound-crossed);
marked partitions exactly equal $(1,1)$. \emph{Controls:} aggregation
(never-coupled) has $\Rrec = \Kbw = (2)$ from the start. \emph{Grade:}
theorem-grade finite instantiation. \emph{Nonclaims:} no phenomenal-unity
derivation; no attention cost or budget (D3 is recognition-grade).

\paragraph{T8 (specious present).}
\emph{Setup:} three-phase clock; mechanism advance $= b/(b+1)$; protocol
permutations on $\{0,1,2\}$. \emph{Result:} funded triple with monotone
arrow; stalls unfunded; noncommuting protocols block a global now; holonomy
alone never certifies flow. \emph{Witness:} progress $2/3$, tick $2/9$, two
wraps, $\Delta A = 2$; holonomy defect $1$ (rotation vs swap).
\emph{Controls:} unfunded $0/0$, constant trace; commuting pair
($R$, $R^2$) defect $0$; flow-null products $3/1000$ each way, blocked.
\emph{Grade:} theorem-grade finite instantiation. \emph{Nonclaims:} no
universal temporal window; no physical time; no spacetime.

\paragraph{T9 (report theorem).}
\emph{Setup:} hidden support $\{u_0, u_1, u_2\}$; report channel; separately
scoped no-report instrument; vendored claim classifier. \emph{Result:}
reports license descended boundary content only; the report map is
non-injective; a no-report channel separates the identified pair.
\emph{Witness:} $(u_0, u_1)$ share $r_\mathrm{shared}$; no-report values
$n_\mathrm{left} \neq n_\mathrm{right}$; overread claim blocked.
\emph{Controls:} visible descent-only claim accepted. \emph{Grade:} direct
instantiation plus finite witness. \emph{Nonclaims:} reports not useless;
hidden content not unreachable by all instruments.

\paragraph{T10 (existence boundary).}
\emph{Setup:} three-state maintenance carrier, requirements $1/2, 3/4, 1$;
record credit (accrue $1/8$ to cap $1/4$; decay $1/8$); budget grid
$\{0, 1/4, 1/2, 3/4, 1\}$. \emph{Result:} graded, hysteretic,
clause-decomposable kernel collapse. \emph{Witness:} up $[0,0,1,2,3]$; down
$[3,3,2,1,0]$; thresholds $1/2$ up vs $1/4$ down; clause patterns
(C1 with collapsed C7; C1 with collapsed C6; full collapse).
\emph{Controls:} weaker rule ($1/16$, cap $1/8$) loses strict hysteresis;
stronger rule ($1/4$, cap $1/2$) never re-crosses on the grid.
\emph{Grade:} schema plus exact correspondences. \emph{Nonclaims:} labels
are structural readings; no clinical claim; no universal hysteresis. Lean
note: sweep constants pinned; island naming and per-point budgets are
Python-only facts.

\paragraph{T11 (valence).}
\emph{Setup:} exact rational LP, maximize $10q_r + 6q_e$ under two budgets.
\emph{Result:} exact binding duals, slack zero, monotone tightening, proxy
instability, viability tie. \emph{Witness:} optimum $q = (4/11, 6/11, 1/11)$,
value $76/11$; $\lambda = (24/11, 14/11)$; curve $(0, 4/3, 24/11, 24/11)$.
\emph{Controls:} slack regime $b = (5,1)$, $\lambda_1 = 0$; proxy phases give
duals $4$ and $0$; planted dual on a slack constraint rejected by
complementary slackness; kernel empties under $1/11$ underfunding.
\emph{Grade:} theorem-grade duality core; the felt-valence identification is
\gradeR{} and retractable. \emph{Nonclaims:} standard duality, no novelty
claim; no generalization to actual subjects.

\paragraph{T12 (quality geometry).}
\emph{Setup:} 32-state micro-carrier (8 hues $\times$ 2 chroma $\times$ 2
texture); lens forgets texture; texture factor $1/2$ breaks lumpability;
staging $\tau = 2$; $\hat P = U P^2 C$, $E = P^2 C U$. \emph{Result:} exact
max-product similarity with full coherence audit; regime-dependent signature
transport. \emph{Witness:} $D = 43/270$ (warm) vs $1/810$ (cool); defect
$13441/52488$; stability max $2/3$, mean $13801/25920$; residuals $180/169$
and $102/59$; signature residue $0$ plateau, $1$ boundary. \emph{Controls:}
swapped prototypes: defect $2619989/3499200$; texture-0 point lift: defect
$167203/524880$, stability max $18/25$, rejected despite $UC = I$.
\emph{Grade:} theorem-grade induced-kernel audit; geometry-birth audit-only.
\emph{Nonclaims:} no curvature construction; no threshold-certified
stability; layer-relative geometry; descended structure only.

\paragraph{T13 (gauge inversion).}
\emph{Setup:} $|Q| = 4$ public classes, two fiber points each; descents:
projection, report, binary fiber path costs; rigidifier: one declared
discrimination-density orientation calibrated by T12's landed signature.
\emph{Result:} gauge group $\mathbb{Z}_2 \to$ trivial under one descended
asymmetry; transports $2 \to 1$. \emph{Witness:} $40320$ automorphisms; $16$
projection-fixing; $G = \{\mathrm{id}, \mathrm{swap}\}$ then
$\{\mathrm{id}\}$. \emph{Controls:} both symmetric gauge elements fix the
public quotient (layer sameness). \emph{Grade:} theorem-grade exact
construction; rigidification principle schema-grade. \emph{Nonclaims:}
rigidifier modeled, not measured; no claim any actual quality space has
nontrivial $G$.

\paragraph{T14 (lineage).}
\emph{Setup:} memory-wheel route package; three histories; $Q = \{q_\mathrm{mid}\}$,
$M = \{C_0, C_1\}$. \emph{Result:} $M$-lineage as identity carrier;
trauma/control loop contrast; fission/fusion classification; empty target.
\emph{Witness:} finer summary factors to $M$, reverse rejected; swap loop
moves $0/1$ of $Q$, $1/1$ of $M$; splitting and merging continuations;
funded kernel size $3$, underfunded $\emptyset$. \emph{Controls:} identity
loop moves nothing; flat/flattenable/currentized comparators refused
lineage. \emph{Grade:} theorem-grade transfer plus finite extensions;
death/token readings deferred \gradeO. \emph{Nonclaims:} no metaphysics of
personal identity; empty-target fact is Python-certified, outside the Lean
kernel.

\paragraph{T15 (combination gate).}
\emph{Setup:} one joint carrier, live states diagonal/crossed plus dead
exit; two interfaces on the same states. \emph{Result:} joint subjecthood
gated by C7 strict refinement and C3 terminal joint reflexivity; aggregation
is the collapsed regime. \emph{Witness:} $\Kbw = (2)$, $\Rrec = (1,1)$;
refinement $1 \to 2 \to 2$; public kernel \{diagonal\} vs terminal kernel
\{diagonal, crossed\}. \emph{Controls:} uncoupled interface has $\Rrec =
\Kbw$; collective ledger exposes pairing (exposure true), failing C4.
\emph{Grade:} schema with exact one-carrier construction. \emph{Nonclaims:}
no verdict on real collectives; aggregation never suffices; Python executes
the refinement operator (Lean proves stability/equivalence).

\paragraph{T16 (meta-problem).}
\emph{Setup:} composite states $(m, q)$; main package fixes $q$; control
tracks $m$; declared completion convention. \emph{Result:} computed flags
$(\top,\top,\top)$ select the hard-problem report; exposed control selects
derivation-certified. \emph{Witness:} split pair $((m_0,q_0),(m_1,q_0))$;
$m$-moved $1$, $q$-moved $0$; no factor $q \to m$; audit
\texttt{undefined\_circular}; overread blocked, visible claim accepted.
\emph{Controls:} exposed interface: factor $\{q_{m0} \mapsto m_0, q_{m1}
\mapsto m_1\}$, flags $(\bot,\top,\bot)$. \emph{Grade:} composition-grade,
convention-relative. \emph{Nonclaims:} not derived from the four laws alone;
premises extensionally coincide on this carrier; illusionism/mysterianism
not refuted.

\paragraph{T18a/b/d/e and T18c.}
\emph{Setup:} T8's clock imported; enablement lenses on $(content, \phi)$;
entrainment gate; T13's transports. \emph{Results and witnesses:} T18a
strictness pair (equal substrate time, different tick parity); T18b coarse
defects $(1,1,1) \to (0,0,0)$ enriched, partial $(4/7)^3 \to (0)^3$ with the
strictly-coarser parity lens, no-birth control all zeros; T18d rows unfunded
(tick $0$), entrained ($1/1$, $1/1$), decoupled ($1/3$, $1/1$); T18e
transports $2$ symmetric, $1$ rigidified. \emph{Controls:} no-birth; P3
null; unfunded. \emph{Grades:} a/b theorem-grade finite, d taxonomy-grade,
e direct reading; \textbf{T18c demoted} to a schema-grade dependency
observation (explicit temporal tokens found only in T1 among T1/T2/T5/T7).
\emph{Nonclaims:} no spacetime; no subject created or detected (T18b); no
sleep/dream claims (T18d); no universal precondition theorem (T18c).

 \section{T17: schema, registry, and negative attacks}
\label{app:t17}


\subsection{Clause-to-witness registry (T17-P)}

The Python protocol's eight rows import actual computed witness fields from
the landed laws; each current row's verdict is \texttt{witness\_verified},
a witness-catalog verdict and never a realization verdict.

\begin{center}\small
\begin{tabular}{llp{0.52\linewidth}}
\toprule
Clause & Source law & Imported witness field (abridged) \\
\midrule
C1 & T1 & closure-objecthood audit: stored-operator checks; unique
$E_2^{(1/16)}$ fixed point retaining $77/80$ CORE mass \\
C2 & T1/T3 & strict pair plus nonfactorization check \\
C3 & T1 & ten-state kernel terminal quotient $6,7,7$, stable, $= M$
restricted; T6 tower as supporting audit \\
C4 & T2 & connector witness with main exposure false \\
C5 & T5 & provenance-routed credit: self / other / none; DPI null blocked \\
C6 & T8 & funded metrics with monotone-ledger arrow; unfunded and
P3-without-P6 nulls \\
C7 & T7 & coupled witness with aggregation collapse control \\
C8 & T1/T11 & full deadline contrast (empty public kernel, nonempty
$L$/no-deadline kernels) plus funded/underfunded viability contrast \\
\bottomrule
\end{tabular}
\end{center}

\subsection{Schema boundaries worth repeating}

The simplified schema's own graded honesty, beyond
Section~\ref{sec:fullvssimp}: its C3 exit check is one-step only (a genuinely
terminal full-domain bisimulation remains an obligation); its C4 is a finite
static specialization of the imported history/complexity-theoretic visibility
predicate, closed under postprocessing of the declared finite joint tuple on
kernel support; its C5 is an aggregate no-fake-arrow balance control (the
concrete two-component lift with actual provenance routing is T5's content,
imported through the source flags); its C6 ledger annotates trace positions
rather than living in the carrier state; its C8 proves strict shrinkage, with
feasible-action lists held fixed; its coherence hub witness is strictly
stronger than C3 and C4 together, deliberately excluding multi-step
determining distinctions at this grade (recorded by an explicit
false-negative boundary fact); and its gate panel's
\texttt{ctrl}/\texttt{nosmuggle}/\texttt{vis}/\texttt{audit}/admissibility
statuses are asserted or procedural at schema grade, because the package
carries no instrument-level records from which to derive them. Full
source-level bridge admissibility, and any theorem about the full design
package, remain priced future work.

\subsection{Attack corpus (abridged)}

Each row is an executable Lean fact rejecting a concrete adversarial package;
paired predecessor facts record genuine old-pass/new-fail surfaces.

\begin{center}\footnotesize
\begin{tabular}{p{0.32\linewidth}p{0.58\linewidth}}
\toprule
Attack family & Representative blocked facts \\
\midrule
Disconnected/off-hub witnesses & disconnected package; disconnected union
(communication isolation); off-hub collapse \\
Hiddenness spoofing & dead-state hiddenness; joint-decoder exposure; HIDD
direction; enabled-action exposure \\
Quotient/refinement decoys & decoy-action refinement (infeasible-action
split); enabled-action merge; dead-state refinement; observation refinement;
coherence decoy pair; multi-step determining boundary (recorded, not
repaired) \\
Credit manufacturing & same-matrix credit; free external matrix; zero-weight
protocol; duplicate-action symmetric mixture \\
Temporal spoofing & one-off tick; drifting tick \\
Viability spoofing & unpriced (zero-cost) existence; unshrinkable kernel;
duplicate kernel mass (two variants) \\
Public-match spoofing & zero-denominator dynamics; undeclared feasible
action; two-step public mismatch; feasibility divergence; channel
divergence; exit leakage; exit-channel complete observation \\
Gate-panel freeloading & \texttt{gate\_panel\_not\_free} (a C2-positive
package cannot acquire a descent-pass panel from an unrelated closure table) \\
\bottomrule
\end{tabular}
\end{center}

Positive boundary controls accompany the corpus (for example, an infeasible
non-lumpable decoy action does not prevent a valid package from matching
itself). The corpus hardens the schema; it does not inhabit it, and
$\texttt{singleCarrierInstanceExhibited} = \texttt{false}$ remains the
registry's load-bearing negative fact.

 \section{Non-evidential empirical-interface preparation history}
\label{app:empirical}


\emph{Protocol.} The artifact contains an admission protocol for
literature-based phenomenon rows. A row is admissible only with: verified
bibliographic metadata; a checked load-bearing passage; a same-sample pair of
altered \emph{and} preserved observables; a proposed clause mapping frozen
independently of the outcome; and a row-level countercondition whose failure
would defeat the row. Mappings are graded as proxies or shape constraints,
never as measurements of a clause.

\emph{Current ruling.} No inherited candidate satisfies every admission
field. The complete inherited battery therefore carries the uniform status
\texttt{history\_only\_non\_evidential}, contributes zero rows to this paper,
and appears nowhere in Sections~\ref{sec:intro}--\ref{sec:limits} as support
for anything. No external C8/existence-boundary row is admissible under the
current human-only, non-neural, report-level scope.

\emph{Reinstatement gates.} A future version may reinstate a row only after:
full-text verification of the load-bearing passage; confirmation of
same-sample altered/preserved observables in the source; an independently
frozen mapping (prospectively registered); and a declared countercondition.
Candidate-level rows, source-status fields, and the append-only rejection log
live in the artifact
(\texttt{paper/prewriting/adequacy\_row\_ledger.csv},
\texttt{paper/prewriting/rejected\_adequacy\_rows.md},
\texttt{paper/prewriting/adequacy\_protocol.md}); they are audit history and
are deliberately not reproduced here.

 \section{Reproducibility and formalization}
\label{app:repro}


\subsection{Verified local baseline}

At the preparation snapshot (2026-07-14, repository working tree): 191
finite-study tests pass across the nineteen law modules; 8 paper-snapshot
tests pass; the deterministic nineteen-law result snapshot
(\texttt{paper/data/current\_results.json}) is current under its strict JSON
schema with drift checking; the 36-job Lean build succeeds; and the
foundations-import, T17-registry, and citation/claim auditors pass, including
their corruption self-tests (each validator is tested against deliberately
corrupted inputs to confirm it fails closed). This is an exact local
baseline, not a release claim.

\subsection{Commands}

From the repository root:
\begin{verbatim}
python paper/prewriting/validate_package.py   # cross-file claim/citation integrity
python paper/tools/export_results.py --check  # snapshot drift check
pytest paper/tests/test_result_snapshot.py -q # snapshot schema/determinism tests
pytest lab/tests -q                           # 191 finite-study tests
lake build                                    # 36-job Lean build
\end{verbatim}
All paper tables and computed figures read
\texttt{paper/data/current\_results.json}; none parses verdict prose or
independently reruns a law module.

\subsection{Theorem-to-artifact map}

Each law has: a Python module
(\texttt{lab/sixbirds\_consciousness/labs/}), a prediction file, a verdict
file with versioned correction blocks, a pytest suite, and a Lean module
(\texttt{lean/SixBirdsConsciousness/Laws/}) with a named kernel declaration
(e.g.\ \texttt{forcingOfI\_T1\_kernel}, \texttt{privacyWheel\_T2\_kernel},
\dots, \texttt{realization\_T17\_kernel}, \texttt{timeMinting\_T18\_kernel}).
The catalog file \texttt{THEOREMS.md} is the canonical statement source; the
T17 contract is machine-readable in
\texttt{formalization/inventory/t17\_claim\_registry.yml} and audited by
script. Per-law Lean scope boundaries are in Table~\ref{tbl:catalog}; the
recurring pattern is that Lean owns core structural facts and pins exhaustive
enumeration outputs as certified constants matching the Python computation.

\subsection{Trust base and known limitations}

Trusted, inherited: vendored foundations Lean modules and their
opacity-guarded axioms (Appendix~\ref{app:imports}); the Lean toolchain; the
Python interpreter and pytest. Known limitations, reported rather than
resolved: the Lean manifest-level validator currently parses zero
inventory/manifest rows due to a schema mismatch and must not be cited as
catalog-wide mechanization evidence (release blocker); no immutable public
artifact DOI or frozen release tag exists yet; release scaffolding (top-level
README, license, citation file, dependency lock, one-command validation) is
incomplete; continuous integration does not yet run the T17 audits or paper
compilation; and the local audit ran on Python~3.10 while the package
declares 3.11+ (to be repeated on a supported version at release). Prediction
files entered version control in the same commits as implementations and are
not preregistration evidence.

 \section{Claim and overlap ledgers}
\label{app:ledgers}


\subsection{Sentence-level claim ledger}

The artifact file \texttt{paper/prewriting/claim\_ledger.csv} is this paper's
prose ceiling: 60 rows, one per allowed claim, each with tier, grade,
allowed-claim wording, evidence, assumptions/scope, the \emph{blocked}
stronger claim, and manuscript destination. Every scientific sentence in this
paper is written to be assignable to one row; the blocked-claim column is the
enforcement surface. The full ledger is reproduced in the artifact rather
than here; its structure is summarized below.

\begin{center}\small
\begin{tabular}{lll}
\toprule
Tier & Rows & Statuses \\
\midrule
S (structural, incl.\ the artifact-audit row) & 50 & allowed (44 overall); allowed\_conditional; allowed\_with\_limitation; demoted \\
R (recognition) & 4 & recognition\_only; future\_protocol\_only \\
O (ontological) & 5 & conditional\_only; assumption \\
OPEN & 1 & open (T17-I) \\
\bottomrule
\end{tabular}
\end{center}

\subsection{Adjacent-paper ownership}

Three adjacent Six Birds papers own seams this paper does not cross
(Section~\ref{sec:related}); the artifact's
\texttt{paper/prewriting/overlap\_memo.md} and
\texttt{sibling\_import\_map.csv} pin every retained sibling object to its
governing version, hypothesis scope, and ownership boundary. The governing
rule: no retained sibling proposition is a logical premise of any theorem
here; sibling results are ownership boundaries or theorem-level analogies,
and ``the same theorem'' is never claimed without a transport proof (none
exists).

\subsection{Disclosures}

\emph{Review status.} This is a preprint and has not been peer reviewed. The
artifact's ``external review request'' documents are author-generated request
artifacts recording an internal review loop.

\emph{Assistance.} Manuscript preparation and repository tooling used AI
assistance; the artifact records the assistance history in
\texttt{paper/prewriting/ai\_assistance\_ledger.md}. The author is
responsible for all claims.

\emph{Funding and competing interests.} To be confirmed by the author before
submission (open item D-01 in the artifact's decision register).
 
\clearpage
\bibliographystyle{plainnat}
\bibliography{prewriting/references_core}

\end{document}
